\section{Bounded, publicly verifiable bridging --- replacing \Fusion, chain consolidation, and why not atomic swaps}
\label{sec:bridging}\label{sec:bridge}

A bridge lets FLUX that lives on the \Flux chain be used on Ethereum or BNB Smart Chain. Today that
works because a single server holds one private key per chain and sends tokens out of a hot wallet
when it observes an inbound payment. The replacement specified in this section locks FLUX in a
publicly visible multi-key reserve on the \Flux L1 and lets one contract per foreign chain issue
exactly that much FLUX there --- never more, because the contract itself refuses. Anyone with a
\Flux node and a public RPC endpoint per chain would be able to check, at any moment, that the
tokens in circulation on every chain add up to what is locked at home. The governing constraint is
the team's own \srcintent{03-parallel-assets.md}: decentralization is the requirement, ``but done
right'' --- \emph{a decentralized design that is exploitable is worse than the custodial
incumbent}. This section takes that literally. It does not attempt to remove every trusted party.
It makes every failure mode bounded by a contract, observable by anyone, and attributable to named
signers, and it states plainly which requirements cannot be met without a \Flux L1 consensus
change.

\medskip
\noindent\textbf{What is now decided, and the numbers that follow.} Six quantities that earlier
drafts of this section carried as open parameters are fixed \srcintent{08-answers-round-2.md, rounds
5--7, 13}: a \emph{global} supply cap of \num{560000000}\flux{} across all chains, enforced by
the reserve invariant under honest-minority operation rather than by any contract (\Cref{rem:global-cap}); a mint rate limit of \num{10000000}\flux per
\SI{24}{\hour} \emph{per chain}, with a matching burn ceiling of \num{10000000}\flux per chain per
window; a 3-of-4 quorum for mint, release and administration alike, held by the four named
co-founders; a guardian pause executable by \emph{any one} of four, with a bounded pause window; no
premine on any surviving chain, the contracts being mint-and-burn so that supply on a chain is only
what has been minted against locked reserve; and limits that are fixed at deployment and not
adjustable --- against contracts that are nonetheless \emph{upgradable} under the same 3-of-4.
Three consequences follow, and this section states all three rather than leaving them to the reader.

\emph{First, the cap is a long-horizon backstop and the rate limit is the bound.}
\num{560000000}\flux is where the emission curve arrives in the early 2050s
(\Cref{sec:monetary-policy}), against a present measured liability of \num{223027781.21}\flux across
all ten legacy assets \srcmeasured{scripts/\allowbreak{}13-issuer-balances.sh}{2026-09-03};
\srcmeasured{scripts/\allowbreak{}15-nonevm-liability.sh, scripts/\allowbreak{}16-remaining-liability.sh}{2026-09-04}; reaching the cap on one
chain at the permitted rate takes \SI{25}{\day}. So the security claim rests on the rate limit, and
because the limit is per chain the worst case scales with the number of live chains: the sustained
figure is \num{20000000}\flux at a \SI{24}{\hour} reaction time over the launch set
$\{\mathrm{ETH}, \mathrm{BSC}\}$ and \num{30000000}\flux once Solana is added
(\Cref{prop:worst-case}). $T_{\mathrm{react}}$, the time to detect a compromise and throw the pause,
is therefore the single most important operational parameter in the bridge.

\emph{Second, one quorum for mint, release and administration collapses the separation the
construction is built around.} The design distinguishes those three roles precisely so that no
single set can both mint and release; with 3-of-4 everywhere, a compromise of three keys obtains
every one of them at once and the rate limit is the only control that still binds. What the 1-of-4
pause preserves is exactly the thing that makes $T_{\mathrm{react}}$ finite: the fourth key can stop
issuance without the other three (\Cref{cor:pause}). That is a single trust assumption bounded by a
rate limit and a pause, and \Cref{sec:bridging-assumption} presents it as such rather than as a
separation of powers.

\emph{Third, and this was the load-bearing open item in the design until $\Delta_A$ was decided
(\Cref{rem:timelock-decided}): an upgradable contract can remove its own rate limit.} A quorum that has been compromised is not obliged to respect a limit it can
delete, and step one of a competent attack is to upgrade the contract rather than to mint under it.
The bound of \Cref{prop:worst-case} therefore holds only if the upgrade path is slower than the
reaction time --- formally, only under the timelock condition $\Delta_A > T_{\mathrm{react}}$ of
\Cref{rem:upgrade-voids}. Without that condition this section can state the cap and cannot state a
worst-case bound at all. \planned

\medskip
\noindent\textbf{Notation used in this section.} \texttt{notation.tex} originally carried no bridge symbols.
The symbols below were local to \Cref{sec:bridging} and were proposed as additions to the canonical macro
file. Three deliberately diverge from the internal mechanism specification to avoid collisions with
symbols that are already canonical: reserve slack is $\varsigma$ (because $\sigma$ is the block
subsidy \subsidy), the mint rate limit is $\varrho$ (because $\rho$ is the \PNR operator
share \nodeshare), and cumulative bridge fees are $F_\bullet$ (because $\Phi$ is the \PNR fee pool
\pool). Cumulative lock and release value are $V_{\mathrm{lock}}$ and $V_{\mathrm{rel}}$ (because
$L$ is the default application lifetime \blockslasting). Existing macros are used wherever they
apply: \hgt, \supply, \flux, \sat, \adv, \advbudget, \defeq, \maxmoney.
These symbols have since been merged into \texttt{notation.tex} as
\texttt{\textbackslash bslack}, \texttt{\textbackslash brate}, \texttt{\textbackslash bcap},
\texttt{\textbackslash breserve}, \texttt{\textbackslash bminted}, \texttt{\textbackslash bflow},
\texttt{\textbackslash bvalue} and \texttt{\textbackslash bquorum}, and appear in the notation index
(\Cref{sec:app-notation}). The cap macro is \texttt{\textbackslash bcap}, rendering $\bcap{c}$;
$\kappa$ is the collusion-set size of \Cref{sec:governance} and is not used in this section.

\begin{table}[H]
\centering\footnotesize
\begin{tabular}{@{}llp{5.6cm}@{}}
\toprule
symbol & domain & meaning \\
\midrule
$\mathcal{C}$, $c$ & --- & the set of canonical parallel chains, and a member of it \\
$M_c$ & $\mathbb{Z}_{\ge 0}$ \sat & minted supply of the canonical token on $c$ \\
$\mathcal{S}_R$ & --- & the published set of L1 reserve scripts (transparent P2SH) \\
$\mathsf{R}$, $\mathsf{R}_{\mathrm{hot}}$, $\mathsf{R}_{\mathrm{cold}}$ & $\mathbb{Z}_{\ge 0}$ \sat & reserve balance, and its hot/cold split \\
$\varsigma$ & $\mathbb{Z}$ \sat & reserve slack, $\varsigma \defeq \mathsf{R} - \sum_c M_c$ \\
$\bcap{c}$ & $\mathbb{Z}_{\ge 0}$ \sat & per-chain contract constant, defence-in-depth only; the \num{560000000}\flux{} cap is \emph{global} and enforced by \Cref{thm:reserve-dominance} under honest-minority operation (\Cref{rem:global-cap}) \\
$\varrho_c$, $\varrho^{\mathrm{burn}}_c$, $T_{\mathrm{win}}$ & $\mathbb{Z}_{\ge 0}$ \sat, time & chain $c$'s contract-enforced mint rate limit and burn ceiling, each decided at \num{10000000}\flux per window; the sliding window, decided at \SI{24}{\hour}. Both are per chain: there is no bridge-wide quantity $\varrho$ \\
$D_{L1}$, $D_c$ & blocks / finality & required confirmation depth on the L1 and on $c$ \\
$(n_M, k_M)$, $(n_R, k_R)$ & $\mathbb{Z}_{>0}^2$ & mint-authorisation quorum; L1 release quorum; both decided at $(4,3)$ \\
$(n_A, k_A)$, $\Delta_A$ & $\mathbb{Z}_{>0}^2$, time & administrative quorum, decided at $(4,3)$, and its timelock \\
$\mathcal{G}$, $\Delta_g$ & set, time & guardian set, and the maximum guardian pause \\
$(n_{\mathrm{att}}, k_{\mathrm{att}})$, $\beta_{\mathrm{att}}$ & $\mathbb{Z}_{>0}^2$, \flux & bonded attester set and per-attester bond \\
$T_{\mathrm{att}}$, $T_{\mathrm{silence}}$, $T_{\mathrm{react}}$ & time & attestation cadence, silence timeout, reaction time \\
$\Lambda_c$, $\Lambda$ & $\mathbb{Z}_{\ge 0}$ \sat & legacy circulating supply on $c$ at the snapshot; $\Lambda = \sum_c \Lambda_c$ \\
$R_0$ & $\mathbb{Z}_{\ge 0}$ \sat & initial L1 backing locked before migration opens \\
$\mathcal{I}_c$ & --- & issuer-held (Foundation) address set on $c$, published before $\hsnap$; its balances are excluded from $\Lambda_c$ \\
$\hsnap$, $h^{\mathrm{mig}}_{\mathrm{close}}$ & height & migration snapshot height; the close of the redemption window, decided at six months from \hpa{} ($\approx \num{3992230}$) \\
$\ell$, $\mathsf{used}_c$ & --- & lock identifier $(\mathrm{txid}, \mathrm{vout})$; consumed lock set on $c$ \\
$n_b$, $B_c$ & $\mathbb{Z}_{\ge 0}$ & a burn nonce; the next burn nonce on $c$ \\
$\varphi_{\mathrm{mint}}$, $\varphi_{\mathrm{rel}}$ & $\mathbb{Z}_{\ge 0}$ \sat & per-operation bridge fees; $F_{\mathrm{mint}}, F_{\mathrm{rel}}$ their cumulative sums \\
$V_{\mathrm{lock}}$, $V_{\mathrm{rel}}$, $V_{\mathrm{top}}$, $V_{\mathrm{out}}$ & $\mathbb{Z}_{\ge 0}$ \sat & cumulative locks, releases paid, other reserve inflows, other outflows \\
\bottomrule
\end{tabular}
\caption{Symbols local to \Cref{sec:bridging}. Every quantity is read at a stated pair of heights
$(\hgt, \{h_c\}_{c \in \mathcal{C}})$.}
\label{tab:bridge-symbols}
\end{table}

\subsection{The incumbent, stated factually}
\label{sec:bridging-incumbent}

Everything in this subsection is \shipped, bar the unfinished rewrite marked \inprogress, and is cited to the deployed source. The system has worked
for years and its operators have recovered from failures by hand. The argument that follows is not
that it has failed; it is that its worst case is unbounded and unobservable, and that a design
exists whose worst case is neither.

\paragraph{What it is.} \Fusion is a Node.js service listening on port \num{9876}
\srccode{fusion/\allowbreak{}config/\allowbreak{}default.js}{6}, backed by MongoDB, that moves FLUX between eleven chains ---
main, KDA, ETH, BSC, TRX, SOL, AVAX, ERG, ALGO, MATIC and BASE
\srccode{fusion/\allowbreak{}config/\allowbreak{}default.js}{106,110}. It has three sub-systems: a snapshot/claim engine that
pays out balances captured by an operator manually running \texttt{flux-\allowbreak{}cli printsnapshot}
\srcdoc{fusion/README.md:27--29}; a swap engine with the state machine
$\textit{new} \to \textit{waiting} \to \textit{confirming} \to \textit{exchanging} \to
\textit{finished}$ and side-states \textit{hold}, \textit{expired}, \textit{dust}
\srccode{fusion/\allowbreak{}src/\allowbreak{}services/\allowbreak{}swapService.js}{861--871}; and a coinbase-claim path. \shipped

\paragraph{Custody.} For each of the eleven chains, \texttt{config/\allowbreak{}default.\allowbreak{}js} holds up to three
role-specific secrets --- \texttt{snapshotsecrets}, \texttt{miningsecrets}, \texttt{swapsecrets} ---
each a single address with a single raw private key, loaded into the process and used to sign
in-process \srccode{fusion/\allowbreak{}config/\allowbreak{}default.js}{342--482}. That is $10 + 11 + 11 = 32$ hot keys, \texttt{main} having no snapshot secret. There
is no MPC, no threshold signature and no multisig anywhere on the path that moves funds out on any
chain \srccode{fusion/\allowbreak{}src/\allowbreak{}coinLibs/\allowbreak{}eth.js}{22,70--101}
\srccode{fusion/\allowbreak{}src/\allowbreak{}coinLibs/\allowbreak{}main.js}{12--22}. The one function whose name contains ``multisig'',
\texttt{signature\allowbreak{}Valid\allowbreak{}Flux\allowbreak{}Multisig}, verifies a \emph{user's} ownership of a P2SH address as a claim
precondition; it says nothing about how outgoing funds are authorised
\srccode{fusion/\allowbreak{}src/\allowbreak{}services/\allowbreak{}generalService.js}{413--432}. What the operator can therefore do
unilaterally is move any balance held under any of the 32 keys to any address, at any time, with no
second party. Whether a cold/hot split exists operationally is indicated only by three named
addresses in a snapshot-exclusion list --- ``Flux Swap Pool LOCKED/HOT/COLD''; no code in the
repository moves funds between them
\srccode{fusion/\allowbreak{}snapshotScriptsMATIC/\allowbreak{}createFluxMainChainSnapshotMATIC.js}{60--70}. \shipped

\paragraph{The in-progress rewrite does not change the model.} In \texttt{fusion-bridge},
\texttt{Bridge\allowbreak{}ZUSDC.\allowbreak{}sol} lets either \texttt{owner} or \texttt{withdrawer} --- plain \texttt{address}
state variables --- call \texttt{release\allowbreak{}USDC(to, amount)} for an arbitrary recipient and amount, with
no proof-of-burn and no Merkle check inside the contract
\srccode{fusion-bridge/\allowbreak{}packages/\allowbreak{}eth-contracts/\allowbreak{}contracts/\allowbreak{}BridgeZUSDC.sol}{13--14,97--102}. The Kadena
side's mint requires a keyset of one public key under a \texttt{keys-all} predicate, i.e.\ 1-of-1
\srccode{fusion-bridge/\allowbreak{}packages/\allowbreak{}kda-contracts/\allowbreak{}src/\allowbreak{}deploy/\allowbreak{}define\_keyset/\allowbreak{}ks.ktpl}{4--7}. The signer
microservice reads its key from a plain environment variable, with Vault integration present only as
commented-out code \srccode{fusion-bridge/\allowbreak{}apps/\allowbreak{}fusion-signer/\allowbreak{}src/\allowbreak{}config/\allowbreak{}keys.ts}{1--2,14--27}.
\inprogress

\paragraph{Correctness of the state machine.} Four facts, each a live code path. \shipped

\begin{enumerate}
\item A swap is written to the database as \texttt{status:\allowbreak{} 'finished'} with
      \texttt{txidTo: 'Sent.\allowbreak{} Awaiting blockchain confirmation.\allowbreak{}.\allowbreak{}.\allowbreak{}'} \emph{before} the outbound send is
      attempted \srccode{fusion/\allowbreak{}src/\allowbreak{}services/\allowbreak{}swapService.js}{893--902}; the real txid is written back
      only under the guard \texttt{if (txid \&\& typeof txid === 'string')}
      \srccode{fusion/\allowbreak{}src/\allowbreak{}services/\allowbreak{}swapService.js}{997--1006}; and every \texttt{send<Chain>()}
      returns the \texttt{Error} object itself on failure
      \srccode{fusion/\allowbreak{}src/\allowbreak{}coinLibs/\allowbreak{}eth.js}{106--112}. A failed send therefore leaves a
      \textit{finished} record carrying a placeholder txid, with no retry, no alert and no rollback.
\item Concurrency guards are module-level booleans (\texttt{swapNotSafe},
      \texttt{some\allowbreak{}Claim\allowbreak{}In\allowbreak{}Process<CHAIN>}, \texttt{automation\allowbreak{}In\allowbreak{}Progress}) and an in-memory
      \texttt{Set}; none survives a restart or a second process instance
      \srccode{fusion/\allowbreak{}src/\allowbreak{}services/\allowbreak{}swapService.js}{443--461}
      \srccode{fusion/\allowbreak{}src/\allowbreak{}services/\allowbreak{}claimSnapshotService.js}{181--190}. The swap-creation guard is a
      busy-wait that proceeds anyway after \SI{8}{\second}
      \srccode{fusion/\allowbreak{}src/\allowbreak{}services/\allowbreak{}swapService.js}{454--458}.
\item Three different confirmation-multiplier tables exist for nominally the same ``enough
      confirmations'' decision \srccode{fusion/\allowbreak{}src/\allowbreak{}services/\allowbreak{}swapService.js}{1085--1090,1181--1219}
      \srccode{fusion/\allowbreak{}src/\allowbreak{}services/\allowbreak{}swapAutomationService.js}{164--182}.
\item Recovery is manual: roughly fifty chain-pair-specific rollback scripts
      (\texttt{helpers/\allowbreak{}adjusToPreviousClaimState*}, \texttt{unholdSwap.js},
      \texttt{remove\allowbreak{}Swap\allowbreak{}TXID.\allowbreak{}js}) and a \textit{hold} state with no automated exit
      \srccode{fusion/\allowbreak{}src/\allowbreak{}services/\allowbreak{}swapService.js}{1216--1219}
      \srccode{fusion/\allowbreak{}helpers/\allowbreak{}processTransactionManual.js}{4--5}.
\end{enumerate}

\paragraph{Gating that exists, and where it lives.} Swap amounts above \num{100000}\flux
(\num{1000000}\flux for a hardcoded premium list) are refused outright
\srccode{fusion/\allowbreak{}src/\allowbreak{}services/\allowbreak{}swapService.js}{150--157}
\srccode{fusion/\allowbreak{}config/\allowbreak{}default.js}{160--195}; swaps still \textit{new} after 90
confirmations go to \textit{hold} \srccode{fusion/\allowbreak{}config/\allowbreak{}default.js}{111}; the blacklist and the
premium list are source-code literals requiring a redeploy to change
\srccode{fusion/\allowbreak{}config/\allowbreak{}blackList.js}{1--27} \srccode{fusion/\allowbreak{}config/\allowbreak{}premiumList.js}{1--12}. The fee is
\SI{0.2}{\percent} plus a flat per-chain amount \srccode{fusion/\allowbreak{}config/\allowbreak{}default.js}{223--236}.
Rate limiting is IP- and user-agent-based with hardcoded Origin/Referer bypass rules
\srccode{fusion/\allowbreak{}src/\allowbreak{}lib/\allowbreak{}server.js}{29,39,54,65,80,91}. All of this gating is enforced in JavaScript
inside one process, not on any chain. \shipped

\paragraph{Reserve accounting.} The EVM tokens were created by minting \num{440000000}\flux to the
deployer in the constructor, at \texttt{decimals =\allowbreak{} 8}
\srccode{flux-eth-bsc/\allowbreak{}flux.sol}{6--13}. Circulating parallel supply is therefore measured as
\texttt{totalSupply()} minus the sum of reserve wallets, and the two team-maintained
reserve-wallet exclusion lists disagree --- neither is a superset of the other, so any such figure
depends on which list is used (see \Cref{sec:parallel-assets}). \shipped

\paragraph{Terminology.} Three separate systems share the name. This paper fixes the usage, and uses
it consistently from here on:

\begin{itemize}
\item \Fusion is the incumbent custodial bridge backend --- the \texttt{Zel\allowbreak{}Core-\allowbreak{}io/\allowbreak{}fusion} service
      described above, plus the unfinished \texttt{fusion-bridge} rewrite that shares its custody
      model. This is what the construction in \Cref{sec:bridging-construction} would replace.
\item \FusionX is the user-facing exchange and redemption product
      \srcdoc{fusion-exchange-frontend/DISCLAIMER.md:1}, whose backend in the corpus is a
      pass-through proxy to environment-configured third-party URLs
      \srccode{fusion-exchange-backend/\allowbreak{}routes/\allowbreak{}trading/\allowbreak{}exchange.js}{96--100}. Its custody model is
      outside the corpus and this paper makes no claim about it.
\end{itemize}

\noindent They are different things, and the distinction matters for
\Cref{sec:bridging-alternatives}: atomic swaps are a plausible mechanism for \FusionX's product and
are not a mechanism for \Fusion's.
\notestablished{the actual custody model behind \FusionX's production backend host is not in the
evidence corpus. The frontend's own ``non-custodial'' claim is unverified and is not repeated here.}

\subsection{Requirements}
\label{sec:bridging-requirements}

The requirements below are derived from the failure modes of \Cref{sec:bridging-incumbent} plus the
team constraint. They reference the participant and adversary definitions of
\Cref{sec:02-system-model}. \textbf{H} denotes hard, \textbf{D} desirable.

\begin{definition}[Dominance over the incumbent, H0]
\label{def:h0}
A replacement design \emph{dominates} the incumbent if, for every participant compromise the design
admits, the worst-case loss is (i) strictly less than the incumbent's --- the entire balance held
under all 32 hot keys, unbounded --- and (ii) publicly observable within a stated bound
$T_{\mathrm{react}}$. A design that fails H0 at any point is rejected regardless of how decentralized
it is.
\end{definition}

\noindent H0 is the team constraint made checkable, and it is the reason the construction leans on
contract-enforced bounds rather than on the number of parties involved. The remaining hard
requirements:

\begin{description}
\item[H1 --- supply is capped by the contract.] On every $c \in \mathcal{C}$, $M_c \le \overline{M}_c$ is
  enforced by the contract, not by an operator, and $M_c$ is readable by anyone.
\item[H2 --- no mint without a lock; no release without a burn.] Every unit minted on $c$
  corresponds to a specific, confirmed, single-use lock on the L1; every unit released on the L1
  corresponds to a specific, confirmed, single-use burn on some $c$.
\item[H3 --- no single key moves funds.] No operation that increases $M_c$ or decreases $\mathsf{R}$
  is authorisable by fewer than $k \ge 2$ independently held keys, and every such authorisation is
  attributable on-chain to the signers that produced it.
\item[H4 --- bounded failure under quorum compromise.] If $k_M$ mint signers collude, the loss is
  bounded by the contract's rate limit and cap, not by the signers' restraint. If $k_R$ release
  signers collude, the loss is bounded by $\mathsf{R}_{\mathrm{hot}}$. H4 is meaningful only against
  a contract the colluding signers cannot rewrite: an upgrade path reachable by the same $k_M$ keys
  reduces H4 to the signers' restraint, which is what \Cref{rem:upgrade-voids} states and what the
  timelock condition there is for.
\item[H5 --- burn is unconditional.] No role --- admin, guardian or attester --- can prevent a holder
  from burning canonical tokens and thereby publishing a claim on the L1 reserve. Pausability applies
  to mint only.
\item[H6 --- auditability from outside.] The reserve invariant of \Cref{sec:bridging-invariant} is
  verifiable by anyone with a \Flux node and public RPC access to each chain, with no privileged data
  and no trust in any team-maintained wallet list. The migration carry alone takes a published
  issuer-held set $\mathcal{I}_c$ as an input, checkable against the legacy contract
  (Algorithm~\ref{alg:carry}); steady-state verification takes none
  \srcintent{08-answers-round-2.md, round 18}.
\item[H7 --- audit before deploy.] Every deployed contract and program has a published independent
  audit of the exact deployed bytecode; migration and sink contracts are in scope
  \srcintent{03-parallel-assets.md}.
\item[H8 --- migration conserves value.] Each legacy unit is redeemable 1:1 into main-chain FLUX
  \srcroadmap{flux-roadmap-public.md:133}; no legacy unit is redeemed twice; and $R_0 \ge \Lambda$
  holds --- the L1 backing for the whole legacy circulating supply exists --- before any migration
  mint occurs.
\item[H9 --- every failure is a stall, not a silent success.] The incumbent's first failure mode ---
  recording success before it happens --- is forbidden by construction: no off-chain record is
  authoritative; chain state is.
\end{description}

\noindent Desirable, not hard: \textbf{D1} redemption liveness under an uncooperative operator
(stated as desirable because \Cref{sec:bridging-open} shows it cannot be made unconditional);
\textbf{D2} minimal admin surface --- the admin cannot mint, transfer, freeze accounts or change
code, and every admin action that weakens a bound is timelocked and public; \textbf{D3}
immutability --- deployed contracts are not upgradeable in place, replacement is by holder-initiated
migration; \textbf{D4} a path to bonded attestation --- pause authority transferable from named
guardians to the attestation network of \Cref{sec:attested-execution} without redeploying the token;
\textbf{D5} one codebase --- ETH and BSC deployments are byte-identical source with chain-specific
constructor parameters only \srcintent{03-parallel-assets.md}.

\paragraph{Tensions, named rather than resolved silently.} Raising $k/n$ makes theft need more
colluders and censorship need fewer refusers; at the decided $(n,k) = (4,3)$ theft needs 3 signers
and censorship needs 2 \srcinferred. No setting maximises both. The design resolves this by making
the \emph{consequences} asymmetric: theft is bounded by the contract (H4) and censorship is bounded
by time, so $k$ may lean toward theft-resistance --- at the cost, at $n = 4$, that the censorship
side is thin, which is why \Cref{op:liveness} is stated as widely as it is. Second, a rate limit
tight enough to bound attacker damage also blocks an exchange migrating tens of millions of tokens at
once, so the migration path is specified to bypass $\varrho_c$ and to be bounded instead by
$\Lambda_c$ and a closing height --- which makes the migrator the largest single mint authority in
the design and puts it squarely in audit scope. Third, H3 asks that no operation moving funds be
authorisable by fewer than $k \ge 2$ \emph{independently held} keys; the decided configuration meets
the arithmetic ($k = 3$) and not the independence, and \Cref{as:admin} says so rather than letting
the requirement appear satisfied. Fourth, D3 asks that a bug be fixed by deploying again and migrating
rather than by upgrading in place --- the price of an audit whose conclusions stay true, and
\Cref{sec:bridging-incidents} records what an in-place upgrade cost Nomad. \textbf{D3 is not
satisfied by the decided configuration}: the contracts are upgradable under the same 3-of-4
\srcintent{08-answers-round-2.md, round 7}, so the construction below specifies immutability while
the decision retains an upgrade path. \Cref{rem:upgrade-voids} states the consequence for
\Cref{prop:worst-case}; it is not a drafting inconsistency to be smoothed over, because the bound
this section exists to state depends on which of the two is deployed.

\subsection{The construction}
\label{sec:bridging-construction}

Everything in this subsection is \planned. Nothing described here is deployed, and no sentence about
it should be read in the present tense.

\begin{definition}[Canonical parallel asset]
\label{def:canonical}
For a chain $c \in \mathcal{C}$, the \emph{canonical} FLUX token on $c$ is the unique contract or
program that is the sole issuer of FLUX on $c$: it has no pre-mint, no owner mint, no proxy, and
issuance only against a signed authorisation referencing a single-use L1 lock, subject to a cap
$\overline{M}_c$ and a rate limit $\varrho_c$ that the contract enforces itself. Its unit is the satoshi
(\texttt{decimals} $=8$ on every chain, matching the legacy EVM tokens
\srccode{flux-eth-bsc/\allowbreak{}flux.sol}{11--13}), so the supply invariant is an integer identity. The decided
launch set is $\mathcal{C}_0 = \{\mathrm{ETH}, \mathrm{BSC}\}$ with
$\mathcal{C}_1 = \mathcal{C}_0 \cup \{\mathrm{SOL}\}$ following \srcintent{03-parallel-assets.md};
one EVM codebase would be deployed twice, plus one Solana program later. \planned
\end{definition}

\begin{definition}[L1 reserve]
\label{def:reserve}
The reserve is a published set $\mathcal{S}_R = \{s_{\mathrm{hot}}, s_{\mathrm{cold},1}, \dots\}$ of
transparent P2SH scripts on the \Flux L1, published together with their redeem scripts so that any
party can recompute both the addresses and the key sets. Its balance is
$\mathsf{R} \defeq \sum_{o \in \mathrm{UTXO}(\mathcal{S}_R)} \mathrm{value}(o)
= \mathsf{R}_{\mathrm{hot}} + \mathsf{R}_{\mathrm{cold}}$. P2SH is used because it is the only
multi-party spending condition \Flux's inherited script offers; the reserve cannot be shielded, both
because a shielded reserve would defeat H6 and because the relevant shielded path is closed after
height \num{835554} \srccode{fluxd/\allowbreak{}src/\allowbreak{}main.cpp}{1398--1408}. Three L1 ledgers are derivable from a
full node: the \emph{lock ledger} (outputs to $s_{\mathrm{hot}}$ whose transaction carries an
\texttt{FLXB1} memo), the \emph{release ledger} (spends from $\mathcal{S}_R$ carrying an
\texttt{FLXR1} memo), and the \emph{top-up ledger} (all other inflows). No off-chain database is a
source of truth (H9). \planned
\end{definition}

\begin{construction}[Canonical mint-and-burn with attested proof-of-reserve, \planned]
\label{con:bridge}
Lock-and-mint from the L1 to $c$; burn-and-release from $c$ to the L1. Four authorities are
separated by role, as in \Cref{tab:bridge-roles}. The contract state on each $c$ is
\begin{align*}
&M_c,\ \overline{M}_c,\ \varrho_c,\ \varrho^{\mathrm{burn}}_c,\ T_{\mathrm{win}},\ \mathsf{bucket}_c,\ \mathsf{bucket}^{\mathrm{burn}}_c
  &&\text{supply, caps, limits, window, buckets}\\
&\mathsf{used}_c \subseteq \{0,1\}^{256} &&\text{consumed lock identifiers}\\
&\mathsf{rem}_c &&\text{unminted remainder per pending lock}\\
&B_c,\ \mathsf{cumBurn}_c,\ \mathsf{cumMint}_c,\ \mathsf{cumMig}_c &&\text{burn counter, cumulative totals}\\
&\Lambda_c,\ h^{\mathrm{mig}}_{\mathrm{close}} &&\text{migration allowance and closing block}\\
&Q_M,\ k_M,\ Q_A,\ k_A,\ \Delta_A,\ \mathsf{queue} &&\text{quorums, timelock, queued changes}\\
&\mathcal{G},\ \mathsf{pausedUntil} &&\text{guardians; mint pause expiry}\\
&\mathsf{migrationTarget},\ \mathsf{predecessor} &&\text{set at most once; fixed at construction}
\end{align*}
with no \texttt{owner}, no \texttt{withdrawer} and no proxy admin slot. The five messages are:

\smallskip
\noindent\textbf{Lock} (L1 transaction). Output $o_1$ pays $v$ \sat to $s_{\mathrm{hot}}$; output
$o_2$ is an \texttt{OP\_RETURN} carrying
$\texttt{"FLXB1"} \| \mathrm{ver}(1) \| \mathrm{chainId}_c(4) \| \mathrm{recipient}(32)$, with
exactly one output to $\mathcal{S}_R$. The lock identifier is
$\ell \defeq (\mathrm{txid}, \mathrm{index}(o_1))$. The transaction is ordinary and needs no
cooperation from any party. The \CumulusVPN entitlement memo \texttt{CVPN1:} is the deployed
precedent for an \texttt{OP\_RETURN}-carried binding on this chain
\srccode{cumulusvpn/\allowbreak{}gateway/\allowbreak{}internal/\allowbreak{}entitle/\allowbreak{}entitle.go}{99--108}.

\smallskip
\noindent\textbf{MintAuth} (off-chain, self-authenticating). For $i \in Q_M$,
\[
\mathrm{sig}_i \;=\; \mathrm{Sign}_{sk_i}\!\Big(H\big(\texttt{"FLXMINT"} \,\|\, \mathrm{chainId}_c
\,\|\, \mathrm{addr}(\mathrm{contract}_c) \,\|\, \ell \,\|\, v \,\|\, \mathrm{recipient}\big)\Big).
\]
The digest is domain-separated by chain identifier \emph{and} contract address, so an authorisation
for one chain cannot be replayed on another or on a successor contract.

\smallskip
\noindent\textbf{Mint} (transaction on $c$). \texttt{mint($\ell$, v, recipient, sigs[])}, submittable
by anyone, admitted by Algorithm~\ref{alg:mint}.

\smallskip
\noindent\textbf{Burn} (transaction on $c$). \texttt{burn(v, flux\allowbreak{}Addr)} emits
$\mathrm{Burn}(n_b, v, \mathrm{fluxAddr}, \mathrm{sender})$ with $n_b \defeq B_c$, then
$B_c \leftarrow B_c + 1$. It is permissionless, has no allowlist, and is never pausable by any role
(H5). On Solana the mint's \texttt{freeze\_authority} is set to \texttt{None} at creation,
irrevocably, for the same reason. The decided configuration adds a burn ceiling
$\varrho^{\mathrm{burn}}_c = \num{10000000}\flux$ per window
\srcintent{08-answers-round-2.md, round 7}, metered by a second leaky bucket
$\mathsf{bucket}^{\mathrm{burn}}_c$ on the same $T_{\mathrm{win}}$. A ceiling on burn is not the same
kind of object as a ceiling on mint --- mint is the direction that can create unbacked value, burn is
the direction that lets a holder leave --- so its interaction with H5 has to be stated rather than
assumed benign, and \Cref{rem:burn-ceiling} does so.

\smallskip
\noindent\textbf{Release} (L1 transaction). Inputs are $k_R$-of-$n_R$ spends of $s_{\mathrm{hot}}$
UTXOs; outputs pay $v_j - \varphi_{\mathrm{rel}}$ to each \texttt{fluxAddr}$_j$ with change to
$s_{\mathrm{hot}}$; an \texttt{OP\_RETURN} carries
$\texttt{"FLXR1"} \| \mathrm{chainId}_c \| n_b^{\mathrm{from}}(8) \| n_b^{\mathrm{to}}(8)$ for a
\emph{contiguous} nonce range. Batches are constructed by Algorithm~\ref{alg:release}.
\end{construction}

\begin{table}[H]
\centering\footnotesize
\begin{tabular}{@{}lp{3.9cm}p{5.4cm}@{}}
\toprule
role & may & may not \\
\midrule
MINT quorum $(n_M,k_M)$ & authorise a mint for a confirmed lock &
  mint beyond $\overline{M}_c$ or $\varrho_c$; mint a used lock; pause; change parameters \\
RELEASE quorum $(n_R,k_R)$ & co-sign L1 spends of $\mathsf{R}_{\mathrm{hot}}$ to burners &
  anything on $c$ \\
GUARDIAN $\mathcal{G}$ (any single member, 1-of-4) & pause mint for $\le \Delta_g$; cancel a
  queued admin action or upgrade \srcintent{08-answers-round-2.md, round 18} &
  pause burn or transfer; mint; release; change parameters \\
ADMIN $(n_A,k_A)$, timelock $\Delta_A$ & change $\overline{M}_c$, $\varrho_c$, quorum sets, $\mathcal{G}$;
  set \texttt{migration\allowbreak{}Target} once; \emph{upgrade the contract}
  \srcintent{08-answers-round-2.md, round 7} & mint; transfer; freeze \\
\bottomrule
\end{tabular}
\caption{The four authorities of \Cref{con:bridge}. The role separation is a property of the
\emph{contract}: the MINT role cannot pause and the ADMIN role cannot mint directly, whoever holds
the keys. It is \emph{not} a property of the signer sets at launch, because the MINT, RELEASE and
ADMIN rows are held by the same 3-of-4 quorum (\Cref{tab:bridge-decided}), so three compromised keys
hold all three at once. The GUARDIAN row is the exception that makes \Cref{prop:worst-case}
evaluable: at 1-of-4 it is reachable by a key the attacker does not hold. The ADMIN row's upgrade
power, however, reaches every other row indirectly --- see \Cref{rem:upgrade-voids}.
\Cref{as:admin} and \Cref{rem:sets-coincide} state the consequence rather than the table implying
otherwise. \planned}
\label{tab:bridge-roles}
\end{table}

\begin{figure}[H]
\centering
\begin{tikzpicture}[
  font=\footnotesize,
  box/.style={draw, rounded corners=2pt, align=center, minimum height=9mm, inner sep=3pt},
  ar/.style={-{Stealth[length=2mm]}, thick}
]
\node[box, minimum width=34mm] (l1) at (0,0) {\Flux L1\\reserve $\mathcal{S}_R$, balance $\mathsf{R}$};
\node[box, minimum width=34mm] (c) at (8,0) {canonical contract on $c$\\supply $M_c \le \overline{M}_c$};
\node[box, minimum width=26mm, dashed] (q) at (4,2.0) {MINT quorum $(n_M,k_M)$\\attributable signatures};
\node[box, minimum width=26mm, dashed] (r) at (4,-2.0) {RELEASE quorum $(n_R,k_R)$\\$k_R$-of-$n_R$ P2SH};
\node[box, minimum width=30mm, dashed] (att) at (8,2.6) {attesters $\to$ guardian\\\texttt{pauseMint()} on $\varsigma<0$};
\draw[ar] (l1) -- node[above, pos=0.5]{Lock $\ell$, depth $D_{L1}$} (q);
\draw[ar] (q) -- node[above, pos=0.5]{Mint, rate-limited} (c);
\draw[ar] (c) -- node[below, pos=0.5]{Burn $n_b$, depth $D_c$} (r);
\draw[ar] (r) -- node[below, pos=0.5]{Release $[n_b^{\mathrm{from}},n_b^{\mathrm{to}}]$} (l1);
\draw[ar, dashed] (att) -- (c);
\end{tikzpicture}
\caption{The two directions of \Cref{con:bridge}. Burn is permissionless and never pausable; mint is
capped, rate-limited and pausable. \planned}
\label{fig:bridge-cycle}
\end{figure}

\paragraph{The decided constants.} \Cref{tab:bridge-decided} fixes the parameters on which every
quantitative claim in this section turns. All nine are decisions, not findings; the provenance is
three intake rounds and the paper carries no independent confirmation of any of them. \planned

\begin{table}[H]
\centering\footnotesize
\begin{tabular}{@{}llp{6.0cm}@{}}
\toprule
parameter & decided value & enforcement point \\
\midrule
$\bcap{c}$ & \num{250000000}\flux & contract constant on each $c$, checked on every mint path; the global \num{560000000}\flux{} cap is enforced by the reserve under honest-minority operation, not here (\Cref{rem:global-cap}) \\
$\varrho_c$ & \num{10000000}\flux \emph{per chain} & leaky bucket inside each contract, Algorithm~\ref{alg:mint} \\
$\varrho^{\mathrm{burn}}_c$ & \num{10000000}\flux \emph{per chain} & second bucket on the burn path (\Cref{rem:burn-ceiling}) \\
$T_{\mathrm{win}}$ & \SI{24}{\hour} & leaky-bucket drain interval for both buckets \\
$(n_M,k_M) = (n_R,k_R) = (n_A,k_A)$ & $(4,3)$ & one signer set for all three, the four named co-founders \\
$\mathcal{G}$ & the same four, \textbf{1-of-4} & any single member pauses mint for $\le \Delta_g$ (\Cref{cor:pause}), or cancels a queued admin action or upgrade \srcintent{08-answers-round-2.md, round 18} \\
limit adjustability & none & $\varrho_c$, $\varrho^{\mathrm{burn}}_c$ fixed at deployment \\
upgrade path & \textbf{present}, under $(4,3)$, behind $\Delta_A = \qty{48}{\hour}$ & admin quorum may replace contract code after a \qty{48}{\hour} public queue (\Cref{rem:upgrade-voids}, \Cref{rem:timelock-decided}) \\
premine & $0$ on every $c$ & no constructor mint; $M_c$ begins at zero \\
\bottomrule
\end{tabular}
\caption{Decided bridge parameters \srcintent{08-answers-round-2.md, rounds 5--7}. The quorum figures
supersede the $(n,k)$ recommendations that \Cref{as:admin} previously carried as open. The guardian
row is no longer marked for confirmation: the pause is decided at 1-of-4, which is what makes
\Cref{prop:worst-case} a bound rather than a restatement of the cap. The last two rows are in
tension with each other --- limits that cannot be adjusted sit inside code that can be replaced ---
and \Cref{rem:upgrade-voids} is where that tension is resolved or not. \planned}
\label{tab:bridge-decided}
\end{table}

\paragraph{Why the limit is per chain and not per address, which is the team's argument and is
correct.}
A per-address mint limit is not a limit. Addresses on every $c \in \mathcal{C}$ are free to create:
an EVM address is the hash of a locally generated public key and costs nothing to produce, so an
adversary subject to a limit of $x$ per address obtains $nx$ by naming $n$ recipients, at the cost of
$n$ transactions and nothing else \srcintent{08-answers-round-2.md, round 5} \srcinferred. Nothing in
the Lock memo or the MintAuth digest of \Cref{con:bridge} binds \texttt{recipient} to an identity, and
no such binding is available on any chain in $\mathcal{C}$ without an identity registry the design
does not have and does not want. A rate control is only a control when the resource it meters cannot
be manufactured by the party being metered; the bucket in Algorithm~\ref{alg:mint} meters issuance
itself, which cannot be. The same reasoning rules out a per-lock or per-transaction ceiling, which an
adversary splits the same way.

\paragraph{Where the limit is enforced, and why a per-chain limit makes that question disappear.}
The contracts are one per chain and cannot read each other's storage, so a limit defined over
$\sum_{c} M_c$ would have no single enforcement point: it could only be a signing policy the mint
quorum observes voluntarily, which bounds nothing in the case the limit exists for, since the party
applying the policy is by hypothesis the compromised party. \textbf{The decision removes that
problem rather than solving it.} $\varrho_c$ is a per-chain constant, and a statement about chain
$c$'s own issuance over $c$'s own clock is exactly what $c$'s contract can check about itself. The
bound therefore holds \emph{under quorum compromise}, because the contract applies it whatever the
signatures say, and no allocation of a shared budget across $\mathcal{C}$ has to be decided,
enforced or revisited. Two consequences are worth naming. Honest throughput is $|\mathcal{C}|
\varrho_c$ rather than a shared $\varrho$, so capacity does not have to be predicted per market ---
which was the operational argument against a fixed split. And the adversary's bound scales the same
way: it is additive in the number of live chains (\Cref{prop:worst-case}), so \emph{adding a chain
raises the worst case by \num{10000000}\flux per reaction window}, which makes chain count a security
parameter and not only a product decision.

\paragraph{No premine, and why that makes the cap unreachable by honest use.} The legacy EVM tokens
were created by minting \num{440000000}\flux to the deployer in the constructor
\srccode{flux-eth-bsc/\allowbreak{}flux.sol}{6--13}. The canonical contracts mint nothing at construction, so
$M_c = 0$ at deployment and every unit on $c$ traces to a specific L1 lock or to a migration deposit
(\Cref{def:canonical}, H2). Combined with \Cref{thm:reserve-dominance} this has a consequence worth
stating: honest issuance is bounded by the reserve, $\sum_c M_c \le \mathsf{R}$, and $\mathsf{R}$
cannot exceed issued main-chain supply, which was \num{429466823.97}\flux at height \num{2912015}
\srcmeasured{emission model, observed mode}{2026-09-03} (\Cref{sec:app-emission}). The cap of
\num{560000000}\flux --- global, and enforced by \Cref{thm:reserve-dominance} rather than by a
contract constant (\Cref{rem:global-cap}) --- therefore cannot bind on any honest path for as long as
issued supply is below it. On the deployed schedule that is the early 2050s: issued supply reaches \num{548043625.86}\flux
at the end of \PoN year 20 (height \num{24095199}, approximately 2046-10) and then grows at a
constant \num{1789219.274256}\flux per year, so \num{560000000}\flux is first issued about
$\num{11956374}/\num{1789219} \approx 6.7$ years later, in 2053 \srcinferred
(\Cref{sec:app-emission}). Only the per-chain constant $\bcap{c}$ binds on the unbacked path, and it
is reached in \SI{25}{\day} at the permitted rate. That is the precise sense in which the cap is a backstop and not
a bound.

\SetKwInput{KwBrInv}{Invariant}
\SetKwInput{KwBrFail}{Failure mode}

\begin{algorithm}[H]
\caption{Contract-side admission of \texttt{mint} on chain $c$. \planned}
\label{alg:mint}
\KwIn{$(\ell, v, \mathrm{recipient}, \mathrm{sigs}[])$; contract state of \Cref{con:bridge}; wall
clock \texttt{now}.}
\KwOut{$v' = v - \varphi_{\mathrm{mint}}$ minted to \texttt{recipient}, or revert.}
\KwBrInv{On any successful return: $M_c \le \overline{M}_c$, $\ell \in \mathsf{used}_c$ once
$\mathsf{rem}_c[\ell] = 0$, $\mathsf{bucket}_c \le \varrho_c$, and $\mathsf{cumMint}_c$ increased by
exactly $u$.}
\KwBrFail{Revert. Execution on $c$ is atomic per transaction, so no partial state change is
possible; the caller may resubmit. No off-chain record is written anywhere (H9).}
\BlankLine
\lIf{$\mathrm{now} < \mathsf{pausedUntil}$}{\Return revert (paused)}
$S \leftarrow$ the set of distinct members of $Q_M$ with a valid signature in $\mathrm{sigs}[]$ over
the MintAuth digest for \emph{this} chain identifier and \emph{this} contract address\;
\lIf{$|S| < k_M$}{\Return revert (quorum)}
\lIf{$\ell \in \mathsf{used}_c$}{\Return revert (lock fully consumed)}
$v' \leftarrow \mathsf{rem}_c[\ell]$ if set, else $v - \varphi_{\mathrm{mint}}$
\tcp*{unminted remainder of a pending lock}
\lIf{$v' \le 0$ \textbf{or} $v < v_{\min}$}{\Return revert (dust)}
$\mathsf{bucket}_c \leftarrow \max\!\big(0,\ \mathsf{bucket}_c - \varrho_c \cdot \Delta t / T_{\mathrm{win}}\big)$
\tcp*{leak since last update}
$u \leftarrow \min\!\big(v',\ \overline{M}_c - M_c,\ \varrho_c - \mathsf{bucket}_c\big)$
\tcp*{admissible now}
\lIf{$u \le 0$}{\Return revert (cap or rate limit; $\ell$ stays pending)}
$\mathsf{rem}_c[\ell] \leftarrow v' - u$;\quad
\lIf{$\mathsf{rem}_c[\ell] = 0$}{$\mathsf{used}_c \leftarrow \mathsf{used}_c \cup \{\ell\}$}
$\mathsf{bucket}_c \mathrel{+}= u$;\quad
$\mathsf{cumMint}_c \mathrel{+}= u$\;
mint $u$ to \texttt{recipient}\;
\end{algorithm}

\paragraph{Why a sliding window and not an epoch counter.} The bucket in
Algorithm~\ref{alg:mint} drains linearly over $T_{\mathrm{win}}$ rather than resetting at fixed epoch
boundaries. With fixed epochs, an adversary who times a compromise at a boundary obtains
$2\varrho_c$ instantly, before any reaction is possible; with a sliding window the total admitted over
an interval of length $T$ is the headroom left in the bucket at the start plus the leak over $T$,
which is at most $\varrho_c(1 + T/T_{\mathrm{win}})$ and is the accounting
\Cref{prop:worst-case} rests on \srcinferred. This is the same class of timing attack that a
deterministic payout schedule admits, and the leaky bucket removes it.

\paragraph{The mint path in full.} A holder broadcasts a Lock. Each mint-quorum signer, from
\emph{its own} L1 node rather than a shared RPC, waits until the lock is at depth $\ge D_{L1}$. The
floor is $D_{L1} \ge 40 = \texttt{MAX\_REORG\_LENGTH}$ \srccode{fluxd/\allowbreak{}src/\allowbreak{}main.h}{74}, because a
reorganisation deeper than 40 blocks is rejected by consensus and the lock is then irreversible for
every node following the rules; at the \PoN target spacing this is
$40 \times \SI{30}{\second} = \SI{20}{\minute}$ \srcinferred. The signer checks memo
well-formedness, that $o_1$ is the only reserve output, and that $\ell \notin \mathsf{used}_c$ read
from its own node on $c$, then publishes its signature. Anyone may then submit the mint; the fee
$\varphi_{\mathrm{mint}}$ was locked and never minted, so it remains in $\mathsf{R}$ as surplus and
requires no fee key. A lock whose value exceeds the bucket's headroom --- including any single lock
above $\varrho_c$ --- is neither refused nor refunded: it stays pending and is minted in instalments
across successive windows, and no admin refund power exists
\srcintent{08-answers-round-2.md, round 18}. Algorithm~\ref{alg:mint} is written for the whole-lock
case; the instalment form keeps an unminted remainder per $\ell$ in place of the binary
$\mathsf{used}_c$ and admits $\min(v',\ \varrho_c - \mathsf{bucket}_c,\ \bcap{c} - M_c)$ per call,
and every bound below is unchanged by it because each instalment passes the same tests.
\settlednote{$D_{L1} = \num{40}$ (\Cref{tab:22-recommended}). $D_{L1} \ge 40$ blocks. Raising it above 40 adds latency and buys nothing against the
2-of-4 emergency-block path \srccode{fluxd/\allowbreak{}src/\allowbreak{}emergencyblock.cpp}{183--191}, which bypasses
eligibility but not the reorg limit. Trade-off: user-visible mint latency against nothing;
40 is both the floor and the chosen value: it is exactly the consensus reorg bound, and exceeding it buys nothing.}
\settlednote{$v_{\min} = \num{10}\flux$, $\varphi_{\mathrm{mint}} = 0$, $\varphi_{\mathrm{rel}} =$ the batch's gas share, recipient Foundation operations (\Cref{tab:22-recommended}). $v_{\min}$ and the fee schedule $(\varphi_{\mathrm{mint}}, \varphi_{\mathrm{rel}})$, plus
the fee recipient. Domain: $v_{\min} > $ L1 dust threshold $+\ \varphi_{\mathrm{rel}}$. Trade-off:
griefing resistance and release-batch size against small-holder access. Constraint: fees must not
scale a signer's income with volume that signer can itself generate.}
\settlednote{A Flux-hosted feed plus a mandatory IPFS mirror (\Cref{tab:22-recommended}). Transport for MintAuth signatures. A Flux-hosted feed is acceptable on the merits ---
signatures are self-authenticating and anyone may relay them, so the feed can withhold but cannot
forge --- but the choice of transport and whether a Flux-independent mirror is mandatory were the
open items, settled as above. Trade-off: operational simplicity against a withholding vector that delays holders who
rely on one relayer.}

\begin{algorithm}[H]
\caption{Release-batch construction, executed independently by each release signer. \planned}
\label{alg:release}
\KwIn{The burn log of chain $c$ up to depth $D_c$; the L1 release ledger for $c$; the UTXO set of
$s_{\mathrm{hot}}$; batch cadence $T_{\mathrm{rel}}$.}
\KwOut{A signature share on a Release transaction covering a contiguous nonce range, or abstention.}
\KwBrInv{The set of nonce ranges released for $c$ is pairwise disjoint and every released nonce is
$< B_c$. Value paid never exceeds $\sum_c \mathsf{cumBurn}_c$ minus what was already released.}
\KwBrFail{Fewer than $k_R$ signers sign $\Rightarrow$ the batch is not broadcast and exits
\emph{stall visibly}; no partial payment occurs and no burn is lost, because the claim remains
on-chain on $c$. A signer that detects an overlapping range abstains, which is the enforcement
point for the disjointness invariant --- the L1 does not check it. A malformed address forfeits its
own burn and blocks no other \srcintent{08-answers-round-2.md, round 18}.}
\BlankLine
$\mathcal{N} \leftarrow$ nonces $n_b$ of burns on $c$ at depth $\ge D_c$, ordered\;
$\mathsf{released}_c \leftarrow$ union of nonce ranges appearing in \texttt{FLXR1} memos for $c$ on the L1\;
$[n^{\mathrm{from}}, n^{\mathrm{to}}] \leftarrow$ the maximal contiguous prefix of
$\mathcal{N} \setminus \mathsf{released}_c$\;
\lIf{$[n^{\mathrm{from}}, n^{\mathrm{to}}] \cap \mathsf{released}_c \ne \emptyset$}{\Return abstain (double
release)}
\lIf{$n^{\mathrm{to}} \ge B_c$}{\Return abstain (nonce beyond counter)}
\ForEach{burn $j$ in the range}{
  verify \texttt{fluxAddr}$_j$ is a well-formed transparent address\;
  \lIf{malformed}{$v_j \leftarrow 0$ (the nonce stays inside the range and is paid nothing; the
  memo records it as \emph{unpayable}, so the prefix advances)}
}
choose inputs $i$ and burners $b$ subject to $34b + 401i + 60 \le \num{2000000}$ bytes
\tcp*{$\le$ \texttt{MAX\_TX\_SIZE\_AFTER\_SAPLING} at the decided 3-of-4}
\Return signature share on the resulting transaction\;
\end{algorithm}

\begin{remark}[The burn ceiling never refuses, and H5 therefore holds --- recommended, not decided]
\label{rem:burn-ceiling}
The burn ceiling $\varrho^{\mathrm{burn}}_c$ admitted two implementations that are indistinguishable
in a parameter table and are not the same object. Under a \emph{refusing} implementation
\texttt{burn} reverts once the window's bucket is full, and a holder is then prevented by a contract
rule from publishing a claim on the reserve --- which is exactly what H5 forbids, and which makes
\Cref{op:liveness} worse rather than better. Under a \emph{delaying} implementation the exit is
deferred to a later window but never refused across the sequence of windows.

\proposednote{A burn is never refused: the delaying implementation, recommended by this paper.
Round 7 decided the \num{10000000}\flux{} ceiling and not which of the two implementations carries
it \srcintent{08-answers-round-2.md, round 7}; no later round chooses, so the choice is recorded here
as a recommendation rather than a decision \srcintent{08-answers-round-2.md, round 18}. Under it H5
holds in the form \emph{no role can prevent a burn from eventually succeeding}, and whether the
implementation carries an explicit deferral queue or simply declines to rate-limit burns at all is
not consequential to any property in this section.}

The reasoning is worth recording, because it also explains why the asymmetry between the two ceilings
is correct rather than an oversight. A mint ceiling protects the reserve: it bounds how much
unbacked supply a compromised quorum can issue, which is the whole content of
\Cref{prop:worst-case}. A burn does the opposite --- it destroys parallel supply and strictly
\emph{improves} the conservation position of \Cref{thm:reserve-dominance}. Limiting it therefore protects
nothing while creating a liveness risk, and the risk lands at the worst possible moment: a burn
ceiling binds during a loss-of-confidence exit, which is precisely the scenario in which holders most
need the exit to work. A refusing ceiling would additionally be cheap to weaponise --- an adversary
could saturate the window with small burns to trap other holders inside it --- so the refusing
reading is not merely worse on principle, it is exploitable.
\end{remark}

\paragraph{The burn path in full.} A holder calls \texttt{burn}; no role can block it. Each release
signer waits for finality on $c$, not a confirmation count: on Ethereum the finalized checkpoint at
two epochs is $2 \times 32 \times \SI{12}{\second} = \SI{768}{\second} = \SI{12.8}{\minute}$, or
\SI{19.2}{\minute} if a third epoch is needed \srcinferred. Signers then co-sign a batch covering a
contiguous nonce range. Contiguity is what makes double-release auditable from the two public
ledgers. Sizing is a bounded, stated quantity rather than an emergent one, and the decided quorum
makes it generous: at $(n_R,k_R) = (4,3)$ the redeem script is $4 \times 34 + 3 = \num{139}$ bytes and
one P2SH input is approximately \num{401} bytes --- three signatures plus the script --- against
\texttt{MAX\_TX\_SIZE\_AFTER\_SAPLING} of \num{2000000} bytes, so one release paying roughly
\num{1000} burners in about \SI{34}{\kilo\byte} of outputs can still consume on the order of
\num{4900} reserve inputs \srcinferred. The same computation at the previously recommended 9-of-15
gave approximately \num{1215} bytes per input and roughly \num{1600} inputs; the small set costs
independence and buys transaction capacity, which is the honest way round to state the trade.
\settlednote{Each chain's own finality tag, never a confirmation count (\Cref{tab:22-recommended}). $D_{\mathrm{ETH}}, D_{\mathrm{BSC}}, D_{\mathrm{SOL}}$ must each be set \emph{at finality},
not at a confirmation count --- the latter is the error class that the BNB Bridge incident belongs
to. Ethereum's value follows from the protocol constant above; BSC's depth under fast finality and
Solana's \texttt{finalized} commitment are deploy-time parameters, set to each chain's own finality predicate rather than to a count. Trade-off:
redemption latency against reorg exposure bounded by $\mathsf{R}_{\mathrm{hot}}$.}
\settlednote{$T_{\mathrm{rel}} = \SI{4}{\hour}$ (\Cref{tab:22-recommended}). Release cadence $T_{\mathrm{rel}}$, domain $[\SI{1}{\hour}, \SI{24}{\hour}]$. Trade-off:
user-visible exit latency against batch efficiency and signer operational load.}
\settlednote{$\mathsf{R}_{\mathrm{hot}} = \num{500000}\flux$, top-up when hot ${<}\ \num{200000}\flux$ (\Cref{tab:22-recommended}). $\mathsf{R}_{\mathrm{hot}}$ target, the cold-reserve threshold, and the cold$\to$hot top-up
cadence. Domain: $\mathsf{R}_{\mathrm{hot}} \approx$ several days of expected redemption volume.
Trade-off: this quantity \emph{is} the L1-side loss bound under H4, against release latency when the
hot balance is exhausted. Sharpened by the quorum decision: the three keys that reach the mint bound
of \Cref{prop:worst-case} are the same three that reach this one, so the exposure of a single 3-key
compromise is $\mathcal{L}(T_{\mathrm{react}}) + \mathsf{R}_{\mathrm{hot}}$ and this target is a
security parameter, not only an operational one.}

\paragraph{Administration, guardianship, and the upgrade path the decision retains.} Any single
$g \in \mathcal{G}$ may call \texttt{pauseMint()}, setting
$\mathsf{pausedUntil} \leftarrow \max(\mathsf{pausedUntil}, \mathrm{now} + \Delta_g)$; the call is
attributable to $g$, and the decided rule is 1-of-4 \srcintent{08-answers-round-2.md, round 7}. The
admin quorum may extend by $\Delta_A$ per signed extension, each on-chain. Burn, transfer and
\texttt{migrate} are never paused, so a pause that is never lifted is a public, attributable,
repeatedly-renewed act by named individuals, which is the whole defence against
pause-as-censorship --- and the bounded window $\Delta_g$ is what keeps a single malicious guardian
able to stall the mint path but not to halt it, since the pause lapses unless renewed and renewal is
again public and attributable. Admin parameter changes are queued on-chain and visible for
$\Delta_A$; \emph{tightening} (lowering $\overline{M}_c$ or $\varrho_c$, removing a signer) executes
immediately, \emph{loosening} waits. Any single $g \in \mathcal{G}$ may also call
\texttt{cancel($q$)} on a queued admin action or upgrade $q$, removing it from $\mathsf{queue}$
before it executes \srcintent{08-answers-round-2.md, round 18}; the call is attributable to $g$, as
the pause is.

\textbf{The construction specifies no proxy, no \texttt{upgradeTo} and no \texttt{delegatecall} to
mutable code (D3); the decided configuration keeps an upgrade path under the same 3-of-4}
\srcintent{08-answers-round-2.md, round 7}. Under D3, replacement on EVM would be by migration:
ADMIN sets \texttt{migration\allowbreak{}Target} once, holders call \texttt{migrate(v)}, which burns on the old
contract and calls \texttt{mintMigrated} on the new one whose \texttt{predecessor} is fixed at
construction; the old contract keeps \texttt{burn} live forever, so a holder who never migrates can
still exit to the L1. The predecessor path is uncapped by $\Lambda_c$ and capped by $\bcap{c}$, and
the burn inside \texttt{migrate} takes no nonce and emits no \texttt{Burn}, so it enters no release
range \srcintent{08-answers-round-2.md, round 18}: were it a nonced burn, Algorithm~\ref{alg:release}
would pay a claim whose value had already been reissued on the successor and $\varsigma$ would go
negative. Under the decided configuration the same replacement is available without
holder action and without notice, which is convenient for repair and is the attack step of
\Cref{rem:upgrade-voids}. The two are not reconcilable by drafting; they are reconcilable only by
the timelock $\Delta_A$, and only if $\Delta_A > T_{\mathrm{react}}$.

Guardian membership at launch is decided: the same four co-founders
\srcintent{08-answers-round-2.md, round 5}, with the single-member pause of \Cref{con:bridge}
retained and now confirmed at source --- \Cref{cor:pause} shows why that retention is load-bearing
rather than incidental. $(n_A, k_A) = (4,3)$ is likewise decided, so the previously open clauses
$n_A \le n_R$ and $k_A \ge k_R$ are satisfied trivially and vacuously: they hold because the sets are
identical, which is not what they were written to secure (\Cref{rem:sets-coincide}).
\settlednote{$\Delta_g = \SI{48}{\hour}$ (\Cref{tab:22-recommended}). $\Delta_g$, domain $[\SI{6}{\hour}, \SI{7}{\day}]$. Note that $\Delta_g > T_{\mathrm{react}}$
is \emph{not} a constraint here: \Cref{prop:worst-case} needs the pause to begin, not to persist, so
what $\Delta_g$ must outlast is only the time to convene the admin quorum and decide what to do.
Trade-off: at 1-of-4 the pause is cheap to trigger by design, so $\Delta_g$ is
exactly the cost a single malicious or mistaken guardian can impose per call. Longer is a stronger
stop and a longer stall; the bound on abuse is that stalling is repeated, public and attributable,
while halting outright is not available at any $\Delta_g$.}
The parameter timelock is likewise $\qty{48}{\hour}$ (\Cref{rem:timelock-decided}), which
satisfies $\Delta_A > T_{\mathrm{react}}$ for any reaction time under two days and leaves a holder
that long to complete a burn-and-release round trip before a parameter change binds them.

\paragraph{Which multisig implementation, and why not the audited one.} On EVM the mint check would
be per-signer ECDSA verification inside the bridge contract --- a distinct-signer bitmap and
\texttt{ecrecover} against $Q_M$ --- and specifically \emph{not} the Halborn-audited Schnorr
\texttt{Multi\allowbreak{}Sig\allowbreak{}Smart\allowbreak{}Account} \srcdoc{account-abstraction/aa-schnorr-multisig-sdk/README.md:76}. Two
evidence-based reasons. First, the aggregated Schnorr signature is verified as a single combined
address, so the contract never learns \emph{which} $k$ signers participated
\srccode{account-abstraction/\allowbreak{}contracts/\allowbreak{}schnorr/\allowbreak{}Schnorr.sol}{7--26}, which defeats H3. Second, its
X-of-Y semantics are realised by enumerating every qualifying signer subset as a separate owner
address \srccode{account-abstraction/\allowbreak{}aa-schnorr-multisig-sdk/\allowbreak{}src/\allowbreak{}helpers/\allowbreak{}schnorr-helpers.ts}{129--152},
i.e.\ $\sum_{j \ge k}\binom{n}{j}$ role grants --- 29 for $(7,5)$ but \num{4944} for $(15,10)$, on the
order of \num{124e6} gas of grants at initialisation \srcinferred --- while the SDK's nonce-reuse
protection is process-local and the documented consequence of reuse is private-key exposure
\srcdoc{account-abstraction/aa-schnorr-multisig-sdk/README.md:14--16}. The Schnorr account is the
right tool for a 2-of-2 wallet and the wrong tool for a bridge quorum. On Solana, the SSP
\texttt{Solana-\allowbreak{}Multisig} program keeps \texttt{approvals:\allowbreak{} Vec<Pubkey>} on-chain and is therefore
attributable \srccode{Solana-Multisig/\allowbreak{}programs/\allowbreak{}solana-multisig/\allowbreak{}src/\allowbreak{}lib.rs}{880--918}, with the
threshold enforced at spend time \srccode{Solana-Multisig/\allowbreak{}programs/\allowbreak{}solana-multisig/\allowbreak{}src/\allowbreak{}lib.rs}{468--471};
it is a candidate for the ADMIN and upgrade authority, but only after an external audit --- none is
recorded in the corpus --- and after its own mainnet upgrade authority is moved off the single
keypair it is on today \srcdoc{Solana-Multisig/docs/MAINNET\_RUNBOOK.md:1--13}. The bridge program's
Solana mint check would verify $k_M$ Ed25519 signatures through the native Ed25519 program with
instruction introspection: precisely the code path whose sysvar check was the Wormhole vector
(\Cref{sec:bridging-incidents}), and therefore a named audit focus.

\noindent\textbf{The quorum question is closed.} $(n_M, k_M) = (n_R, k_R) = (4,3)$, one set, the four
named co-founders \srcintent{08-answers-round-2.md, round 5}. Both sets sit well inside the P2SH
ceiling --- standardness caps the redeem script at 520 bytes, which is 15 compressed keys
\srcinferred --- so the constraint that bounded $n_R$ from above is not active. The theft/censorship
trade-off resolves as follows at $(4,3)$: theft needs 3 colluders and censorship needs 2 refusers, so
the release path is easier to stall than to rob, and \Cref{op:liveness} is correspondingly wider than
it would be at a larger $n_R$. \Cref{sec:bridging-assumption} states what the shared membership costs.
\settlednote{The named hardware class and procedure of \S\ref{sec:22:recommended}. Custody standard for the four keys: dedicated signing hardware, keys never on a host
running a relayer, an RPC or the incumbent service, and a documented key-generation ceremony.
Domain: a named hardware class plus a published procedure. Trade-off: a stricter standard slows
signing and reduces the chance that one compromise of a general-purpose host yields a key --- which
under a single 3-of-4 set yields every capability at once. Note that the clauses about
\emph{organisational} distribution and jurisdictional spread are not open parameters any more; they
are not satisfied at all, and \Cref{as:admin} says so.}

\subsection{The supply invariant}
\label{sec:bridging-invariant}

\begin{theorem}[Reserve dominance]
\label{thm:reserve-dominance}
Let all quantities be read at a stated pair of heights $(\hgt, \{h_c\})$, and assume
(i) every mint on every $c$ passed the checks of Algorithm~\ref{alg:mint};
(ii) at most $k_M - 1$ mint signers and at most $k_R - 1$ release signers are corrupt on each chain;
(iii) $R_0 \ge \Lambda$ (H8); and
(iv) $V_{\mathrm{out}} = 0$, i.e.\ no reserve outflow occurs other than releases.
Then
\begin{align*}
\varsigma \;\defeq\; \mathsf{R} - \sum_{c \in \mathcal{C}} M_c
\;=\;& \underbrace{\Big(\Lambda - \sum_c \mathsf{cumMig}_c\Big)}_{\text{(1) unmigrated legacy}}
\;+\; \underbrace{\Big(V_{\mathrm{lock}} - \sum_c \mathsf{cumMint}_c - F_{\mathrm{mint}}\Big)}_{\text{(2) locks not yet minted}} \\[0.6em]
{}+\;& \underbrace{\Big(\sum_c \mathsf{cumBurn}_c - V_{\mathrm{rel}} - F_{\mathrm{rel}}\Big)}_{\text{(3) burns not yet released}}
\;+\; \underbrace{\Big((R_0 - \Lambda) + V_{\mathrm{top}} + F_{\mathrm{mint}} + F_{\mathrm{rel}} - V_{\mathrm{out}}\Big)}_{\text{(4) surplus}}
\end{align*}
and every bracket is non-negative; hence $\varsigma \ge 0$, i.e.\ $\mathsf{R} \ge \sum_c M_c$, at
every pair of heights. \planned
\end{theorem}

\begin{proof}[Proof sketch]
The contract's arithmetic admits exactly three mutators of \texttt{totalSupply}, so
$M_c = \mathsf{cumMig}_c + \mathsf{cumMint}_c - \mathsf{cumBurn}_c$. UTXO accounting over
$\mathcal{S}_R$ gives $\mathsf{R} = R_0 + V_{\mathrm{lock}} + V_{\mathrm{top}} - V_{\mathrm{rel}}
- V_{\mathrm{out}}$, where $V_{\mathrm{rel}}$ counts only value paid to burners --- fees withheld at
release never leave the reserve and are therefore counted once, in bracket~(4). Substituting and
grouping yields the displayed identity by cancellation of $\Lambda$, $F_{\mathrm{mint}}$ and
$F_{\mathrm{rel}}$.

Non-negativity, bracket by bracket. (1) The migration mint path is bounded by
$\mathsf{cumMig}_c + v \le \Lambda_c$ inside the contract, and $\sum_c \Lambda_c = \Lambda$. (2)
Each public mint consumes a distinct $\ell$ from the lock ledger and mints $v' = v -
\varphi_{\mathrm{mint}} \le v$; single-use is enforced by $\mathsf{used}_c$ under consensus on $c$,
and value correctness by the honest signer that assumption (ii) guarantees is present in any set of
$k_M$ valid signatures. (3) Each release pays a nonce range contained in $[0, B_c)$, and
disjointness is enforced by the abstention rule of Algorithm~\ref{alg:release}, which runs whenever
at least one of the $k_R$ required signers is honest --- again assumption (ii). (4) $R_0 - \Lambda
\ge 0$ by assumption (iii), $V_{\mathrm{top}} \ge 0$ and $F_\bullet \ge 0$ by construction, and
$V_{\mathrm{out}} = 0$ by assumption (iv).

Two dynamic cases. \emph{Across the migration}: bracket~(1) drains into $\sum_c M_c$ while $R_0$
stays in $\mathsf{R}$, so $\varsigma$ decreases monotonically toward the surplus and never below it.
\emph{Across a successor contract}: \texttt{migrate} burns on the predecessor and mints on the
successor within one transaction, so $\sum_c M_c$ is unchanged and $\varsigma$ is invariant; neither
leg is a nonced burn or a migration mint, so brackets~(1) and~(3) are unchanged as well.
\end{proof}

\begin{corollary}[The quiescent form, and what the roadmap says]
\label{cor:quiescent}
Write $\supply_{\mathrm{main}} \defeq \supply(\hgt) - \mathsf{R}$ for main-chain circulating supply.
In the quiescent state --- no operation in flight, migration complete
($\sum_c \mathsf{cumMig}_c = \Lambda$), surplus swept --- every bracket vanishes, so $\varsigma = 0$
and therefore $\sum_c M_c + \supply_{\mathrm{main}} = \supply(\hgt)$. That identity is the roadmap's
phrasing, ``minted supply across every chain plus main-chain supply equals total issuance''
\srcroadmap{flux-roadmap-public.md:99}. The live form is the inequality $\varsigma \ge 0$ with the
decomposition of \Cref{thm:reserve-dominance}, which is strictly more useful: it is what makes a
\emph{negative} $\varsigma$ diagnosable, because each bracket going negative names a distinct
failure --- migrator overrun, unbacked mint, unbacked release, or reserve theft --- and each is
attributable to exactly one role in \Cref{tab:bridge-roles}. \planned
\end{corollary}

\begin{algorithm}[H]
\caption{Public verification of the reserve invariant. Runnable by anyone. \planned}
\label{alg:verify}
\KwIn{A \Flux full node; one RPC endpoint per $c \in \mathcal{C}$; the published $\mathcal{S}_R$ with
redeem scripts; a chosen height pair $(\hgt, \{h_c\})$.}
\KwOut{$\varsigma$ and the four brackets of \Cref{thm:reserve-dominance}.}
\KwBrInv{The computation reads only public chain state. It requires \textbf{no} team-maintained
wallet list and no privileged endpoint.}
\KwBrFail{If any bracket is negative, report it together with the role it implicates. If an RPC is
unreachable, the run aborts rather than substituting a cached value --- a partial read must not
produce a reassuring number.}
\BlankLine
recompute each address in $\mathcal{S}_R$ from its published redeem script; abort on mismatch\;
$\mathsf{R} \leftarrow$ sum of unspent outputs to $\mathcal{S}_R$ at height $\hgt$\;
$V_{\mathrm{lock}} \leftarrow$ sum over the \texttt{FLXB1} lock ledger;\quad
$V_{\mathrm{rel}} \leftarrow$ sum paid to burners over the \texttt{FLXR1} release ledger\;
\ForEach{$c \in \mathcal{C}$}{
  read $M_c, \mathsf{cumMint}_c, \mathsf{cumMig}_c, \mathsf{cumBurn}_c, B_c$ from contract storage at
  $h_c$\;
  check the released nonce ranges for $c$ are pairwise disjoint and all $< B_c$\;
}
$\varsigma \leftarrow \mathsf{R} - \sum_c M_c$; evaluate brackets (1)--(4)\;
\Return $\varsigma$ and the brackets\;
\end{algorithm}

\paragraph{Who can verify, and why the absence of a wallet list matters.} Anyone. $\mathsf{R}$ and
the three L1 ledgers are transparent chain facts by \Cref{def:reserve}; $M_c$ and the four
cumulative counters are public contract storage on $c$. A verifier needs a \Flux full node and one
RPC endpoint per chain, and needs \emph{no} team-maintained wallet list. This is not a stylistic
point. \Cref{sec:parallel-assets} establishes that two team-maintained reserve-wallet exclusion
lists exist today and that they disagree --- neither is a superset of the other --- so any current
circulating-supply figure that nets out team reserves depends on which list the reader happens to
use. Having exactly one published $\mathcal{S}_R$, recomputable from published redeem scripts,
removes that dependency entirely. Total issuance $\supply(\hgt)$ is itself verifiable from
\texttt{gettxoutsetinfo} together with \texttt{valuePools}, with the standing caveat that ZIP-209
turnstile enforcement is compiled in but dormant
\srccode{fluxd/\allowbreak{}src/\allowbreak{}chainparams.cpp}{322--327}; note that \Cref{thm:reserve-dominance} does not
depend on that caveat, because it compares $\mathsf{R}$ against $\sum_c M_c$ and never against
$\supply(\hgt)$.

\begin{proposition}[Bounded loss under full mint-quorum compromise]
\label{prop:bounded-mint}
Drop assumption (ii) of \Cref{thm:reserve-dominance} and suppose $k_M$ or more mint signers on chain
$c$ are corrupt from time $t_0$, with a pause first effective at $t_0 + T_{\mathrm{react}}$. Write
$b_c(t) \in [0, \varrho_c]$ for the bucket fill. Then the unbacked issuance they can produce on $c$
is at most
\[
\big(\varrho_c - b_c(t_0)\big) \;+\; \varrho_c \cdot \frac{T_{\mathrm{react}}}{T_{\mathrm{win}}}
\;\le\; \varrho_c\Big(1 + \frac{T_{\mathrm{react}}}{T_{\mathrm{win}}}\Big),
\]
and at most $\bcap{c} - M_c(t_0)$ in total; every such mint is an on-chain transaction from which the
colluding signers' identities are recoverable. \planned
\end{proposition}

\begin{proof}[Proof sketch]
Algorithm~\ref{alg:mint} applies the cap and bucket checks after the quorum check and regardless of
its outcome, so no set of signatures can exceed them; corruption removes only the guarantee that a
lock exists, never the guarantee that the bounds hold. For the bucket term, every admitted mint of
$v'$ increases $b_c$ by exactly $v'$ and the only decrease is the leak at rate
$\varrho_c / T_{\mathrm{win}}$, so the total admitted over $[t_0, t_0 + T]$ equals
$b_c(t_0 + T) - b_c(t_0) + \varrho_c T / T_{\mathrm{win}}$, and $b_c(t_0+T) \le \varrho_c$ by the
admission test. Mints authorised before $t_0$ and submitted after it consume the same bucket, so they
displace rather than add to the adversary's share; the earlier phrasing ``plus in-flight
authorisations'' overstated the bound and is withdrawn. Attributability follows from per-signer
ECDSA verification against $Q_M$: the signatures are contract inputs and are therefore permanently
public.
\end{proof}

\noindent With the parameters of \Cref{tab:bridge-decided} this becomes a number, and it is the
number on which the section's security claim now rests. Two hypotheses have to be carried
explicitly, because each of them is a decision that can be taken either way and the property is
false without them.

\begin{property}[Worst-case loss of the replacement bridge]
\label{prop:worst-case}
Let $\varrho^\star \defeq \num{10000000}\flux$ be the decided per-chain mint rate limit, so
$\varrho_c = \varrho^\star$ for every $c \in \mathcal{C}$, with
$T_{\mathrm{win}} = \SI{24}{\hour}$ (\Cref{tab:bridge-decided}). Suppose the 3-of-4 quorum of
\Cref{as:admin} is compromised at $t_0$ and the adversary mints as fast as
Algorithm~\ref{alg:mint} admits on every $c$. Assume
\begin{description}
\item[(P)] a pause is first effective at $t_0 + T_{\mathrm{react}}$ with $T_{\mathrm{react}}$ finite,
  which is what the 1-of-4 guardian rule of \Cref{cor:pause} supplies; and
\item[(U)] the adversary cannot replace the contract's code inside $T_{\mathrm{react}}$, which
  requires an upgrade timelock $\Delta_A > T_{\mathrm{react}}$ (\Cref{rem:upgrade-voids}).
\end{description}
Then the unbacked issuance over $[t_0,\, t_0 + T_{\mathrm{react}}]$ is
\[
\mathcal{L}(T_{\mathrm{react}})
\;\le\;
\underbrace{|\mathcal{C}|\,\varrho^\star \cdot \frac{T_{\mathrm{react}}}{T_{\mathrm{win}}}}_{\text{sustained}}
\;+\;
\underbrace{\textstyle\sum_{c \in \mathcal{C}}\big(\varrho^\star - b_c(t_0)\big)}_{\text{headroom at } t_0,\ \le\ |\mathcal{C}|\,\varrho^\star}
\qquad\text{and}\qquad
\mathcal{L} \;\le\; \sum_{c \in \mathcal{C}}\big(\bcap{c} - M_c(t_0)\big).
\]
The sustained term is additive in the number of live chains and evaluates as in
\Cref{tab:worst-case}. The headroom term adds at most one further $|\mathcal{C}|\,\varrho^\star$ ---
\num{20000000}\flux over the launch set --- and is largest exactly when the bridge is idle at the
moment of compromise. The cap term does not bind on any chain for
$T_{\mathrm{react}} < \SI{25}{\day}$. \planned
\end{property}

\begin{proof}
Sum the per-chain bound of \Cref{prop:bounded-mint} over $c \in \mathcal{C}$ with
$\varrho_c = \varrho^\star$ on both the sustained and the headroom term; the cap bound is the check
$M_c + v' \le \bcap{c}$, which Algorithm~\ref{alg:mint} applies on every path. Hypothesis (P) is what
makes $T_{\mathrm{react}}$ a finite quantity to substitute, and \Cref{cor:pause} discharges it.
Hypothesis (U) is what makes Algorithm~\ref{alg:mint} the code that runs for the whole of
$[t_0, t_0 + T_{\mathrm{react}}]$; \Cref{prop:bounded-mint} bounds what \emph{that} algorithm admits
and says nothing about a successor contract, so without (U) the sum is over an empty set of
enforced checks. For the numbers, with $|\mathcal{C}_0| = 2$ at launch:
$2\varrho^\star \cdot (24/24) = \num{20000000}\flux$, $2\varrho^\star \cdot (48/24) =
\num{40000000}\flux$, and $3\varrho^\star \cdot (24/24) = \num{30000000}\flux$ once
$\mathcal{C}_1 = \mathcal{C}_0 \cup \{\mathrm{SOL}\}$; and
$\bcap{c} / \varrho^\star = \num{250000000}/\num{10000000} = 25$ windows of \SI{24}{\hour}, so from
an empty contract the cap on one chain is first reached at \SI{25}{\day}, which exceeds any
$T_{\mathrm{react}}$ the design contemplates \srcinferred.
\end{proof}

\begin{table}[H]
\centering\footnotesize
\begin{tabular}{@{}llr@{}}
\toprule
chains live & $T_{\mathrm{react}}$ & sustained worst-case mint \\
\midrule
$\mathcal{C}_0 = \{\mathrm{ETH}, \mathrm{BSC}\}$ (launch) & \SI{24}{\hour} & \num{20000000}\flux \\
$\mathcal{C}_0 = \{\mathrm{ETH}, \mathrm{BSC}\}$ & \SI{48}{\hour} & \num{40000000}\flux \\
$\mathcal{C}_1 = \mathcal{C}_0 \cup \{\mathrm{SOL}\}$ & \SI{24}{\hour} & \num{30000000}\flux \\
\bottomrule
\end{tabular}
\caption{The sustained term of \Cref{prop:worst-case} at the decided per-chain limit
$\varrho^\star = \num{10000000}\flux$ per \SI{24}{\hour}
\srcintent{08-answers-round-2.md, round 7}. Headroom at $t_0$ adds at most a further
$|\mathcal{C}|\,\varrho^\star$. These figures hold under hypotheses (P) and (U) and not otherwise.
\planned}
\label{tab:worst-case}
\end{table}

\begin{corollary}[The pause is reachable by a party outside any 3-key compromise]
\label{cor:pause}
\Cref{prop:worst-case} is a statement about $T_{\mathrm{react}}$, so it would be vacuous if the
adversary also held the pause. The adversary does not. The decided rule is that \texttt{pauseMint()} is a
\textbf{single-member action of $\mathcal{G}$ at $|\mathcal{G}| = 4$} --- 1-of-4, with the bounded
window $\Delta_g$ \srcintent{08-answers-round-2.md, round 7}. An adversary holding three of the four
keys therefore does not control the pause: one guardian remains able to act, and
$T_{\mathrm{react}}$ is bounded above by that member's detection-and-response time. Hypothesis (P)
of \Cref{prop:worst-case} is discharged, and the worst-case bound holds rather than collapsing to
the cap. \planned
\end{corollary}

\begin{proof}[Proof sketch]
$T_{\mathrm{react}}$ enters \Cref{prop:worst-case} only through the time at which
$\mathsf{pausedUntil}$ first exceeds \texttt{now} in Algorithm~\ref{alg:mint}. That state is written
only by \texttt{pauseMint()}, whose authorisation predicate is the guardian rule. Under the decided
single-member rule the predicate is satisfiable by any one of four keys; an adversary holding three
cannot prevent the fourth from broadcasting, because the call is a permissionless transaction on $c$
once signed, and nothing in Algorithm~\ref{alg:mint} lets a mint consume or reset
$\mathsf{pausedUntil}$. The bound is on $T_{\mathrm{react}}$ as a \emph{human} latency, not a
protocol one: the honest guardian must notice, decide and sign, which is the subject of the
paragraph after \Cref{rem:upgrade-voids}. Note what the proof does \emph{not} establish --- that the
pause holds. It expires at $\Delta_g$ unless renewed, and renewal against a compromised admin quorum
is a repeated race the guardian wins only by repeating the call.
\end{proof}

\paragraph{Why the pause is cheap and the mint is not, which is the design reason and not a
compromise.} Pausing and minting have opposite requirements. \textbf{Pausing is a safety action that
should be cheap to trigger and expensive to abuse; minting is a value-creating action that should be
expensive to trigger and is not abusable at all if the trigger holds.} A high threshold on the mint
quorum buys exactly what a high threshold is for --- more colluders needed to create value from
nothing. The same threshold on the pause buys nothing and costs the bound, because the set that must
be unable to mint is the set that would then also be able to refuse to stop. So the two thresholds
are set in opposite directions, 3-of-4 and 1-of-4, and the asymmetry is the point rather than an
inconsistency. What makes the cheap trigger tolerable is the bounded window: a single malicious or
mistaken guardian can impose $\Delta_g$ of stalled minting per call, publicly and attributably, and
cannot halt the bridge, because burn, transfer and \texttt{migrate} are never pausable (H5) and the
pause lapses on its own.

\paragraph{Why the bound is proved against quorum compromise at all, since the objection is fair.}
A reader may reasonably ask why the section spends its central property on a scenario in which three
of four named co-founders' keys are in an attacker's hands --- a bar that is high, and this section
says so. Two answers, and the section states them rather than assuming the reader shares the
premise. First, \textbf{a bound that holds only when nobody is compromised is not a security
property.} It is a description of normal operation, which the contract's ordinary behaviour already
supplies. The entire value of a rate limit is what it guarantees in the case one hopes never happens;
evaluated against an honest quorum it guarantees nothing that was in doubt. Second, \textbf{the case
is not hypothetical in this domain.} Ronin was a 5-of-9 validator compromise, and
\Cref{tab:bridge-incidents} lists it alongside other multi-signer quorums that were taken whole. The
bar being high is a real and stateable mitigation, and it is the reason the design is defensible at
launch (\Cref{sec:bridging-assumption}); it is not a reason to leave the case unmodelled. \planned

\begin{remark}[Upgradability voids the rate limit unless the upgrade is slower than the reaction]
\label{rem:upgrade-voids}
The limits are fixed at deployment and not adjustable --- and the contracts are upgradable under the
same 3-of-4 \srcintent{08-answers-round-2.md, round 7}. An adversary holding three keys is not
obliged to mint under a limit it can delete: \textbf{step one of the attack is to upgrade the
contract and remove the limit}, after which the bound of \Cref{prop:worst-case} constrains code that
is no longer running. As specified today the rate limit therefore bounds nothing under exactly the
scenario it exists to bound, and the strongest true statement available is the cap term,
$\sum_c(\bcap{c} - M_c(t_0))$ --- a quantity starting from $\bcap{c} = \num{250000000}\flux$ per
chain against a present outstanding liability of \num{223027781.21}\flux across all ten measured
legacy chains \srcmeasured{scripts/\allowbreak{}13-issuer-balances.sh}{2026-09-03}, and therefore more than an
order of magnitude weaker than \Cref{tab:worst-case}. Note that the 1-of-4 pause does not rescue
this on its own: an upgrade that removes the limit can remove the guardian rule in the same
transaction.

\textbf{The fix composes with what is already decided and costs nothing new.} Put the upgrade path
behind the admin timelock $\Delta_A$ and require
\[
\Delta_A \;>\; T_{\mathrm{react}},
\]
which is hypothesis (U). Then a malicious upgrade is queued on-chain and publicly visible before it
takes effect; the 1-of-4 pause can fire during the delay; the limit is still the code that runs for
the whole reaction window; and the rate limit binds again. The mechanism is the one the construction
already specifies for loosening parameter changes, applied to the code as well as to the constants.
Cancellation of a queued upgrade or admin change is likewise a single-member guardian action
\srcintent{08-answers-round-2.md, round 18}; without it the pause would defer the upgrade's effect
only until $\Delta_g$ lapses, the queued upgrade would execute at $t_0 + \Delta_A$ regardless, and
the bound of \Cref{prop:worst-case} would end there with $\sum_c \bcap{c}$ as the only ceiling
after it.
\textbf{Without a timelock on the upgrade path, this section can state the cap and cannot state a
worst-case bound at all}, and every figure in \Cref{tab:worst-case} is conditional on it. This was
the load-bearing open item in the bridge design; the decision follows. \planned
\begin{remark}[$\Delta_A$ decided: \qty{48}{\hour}]
\label{rem:timelock-decided}
The upgrade timelock is set at $\Delta_A = \qty{48}{\hour}$, with upgrades remaining $3$-of-$4$
\srcintent{evidence/design/08-answers-round-2.md}. Unanimity was considered and rejected for a
reason worth recording, because it is the opposite of the one that first suggests itself: $4$-of-$4$
would indeed prevent an adversary holding three keys from upgrading, but it has no fault tolerance
at all --- a single lost, destroyed or unavailable key would make the contract permanently
unupgradable, converting a key-theft risk into a key-loss risk, and for signer sets this small key
loss is the more common failure. The recovery path from that state is a fresh deployment and another
migration, which is precisely the event upgradability exists to avoid.

\textbf{Upgradability is a bring-up mechanism with immutability as its stated destination.} The
team's reasoning for keeping it, and for keeping the quorum at 3-of-4, is that a management path is a
safety net while the system is being verified to work without bugs, and that ``it's easier to upgrade
contract rather than releasing a new one and doing migration''; as the network evolves the intent is
to move to non-upgradable contracts and automate what the quorum does today
\srcintent{08-answers-round-2.md, round 11} (\vision). The paper records this because it changes what
the trust assumption \emph{is}: not a permanent property of the design, but an explicitly temporary
one whose removal is intended and undated.

That framing is coherent, and the honest statement of its weakness is a matter of sequencing rather
than of principle. An upgradable contract under a 3-of-4 quorum is, for as long as it lasts, a
contract whose rules those three signers can rewrite --- the \SI{48}{\hour} timelock bounds the speed
of that rewrite, not its scope (\S\ref{sec:bridging-assumption}). A commitment to eventual immutability
constrains nothing until it has a trigger, and none is stated: no value threshold, no elapsed-time
condition, no audit milestone. \Cref{tab:22-recommended} does not assign one, because a date the
team has not chosen would be an invention. What the paper can say is what a trigger would need to
look like --- a published condition, checkable by a third party, after which the upgrade path is
renounced on chain --- and that until such a condition exists, the bridge should be described as
administered rather than trust-minimized, which is how \S\ref{sec:bridging-assumption} describes it.
\paragraph{The renunciation trigger: stated in shape, proposed here in checkable form.} Two
preconditions for locking immutability are settled: a \textbf{third-party security audit}, and a
\textbf{proven operating record across edge scenarios --- chain reorganizations named explicitly}
\srcintent{08-answers-round-2.md, round 12} (\planned). Separately, the team states that code will
primarily be written and audited using state-of-the-art AI models, with the third-party audit as an
additional and distinct requirement rather than a substitute
\srcintent{08-answers-round-2.md, round 12}. The paper records that as the stated development
practice; it is a claim about process, and this document takes no position on the assurance it
provides, which is why the separate third-party requirement is the load-bearing one here.

Both preconditions are the right ones, and neither is yet a trigger, because a trigger has to be
checkable by someone who does not work on the project. ``A proven time record'' is not falsifiable
without a duration, and --- more importantly --- \textbf{elapsed time proves nothing on a bridge
nobody uses}. A contract that sits idle for a year has demonstrated nothing about its behaviour
under load, under reorganization, or under an adversary; it has demonstrated only that it was idle.
Any time-based condition therefore needs a volume condition beside it, or it can be satisfied by
inactivity. The following formulation is this paper's proposal, chosen so that every clause is
verifiable from public data:

\begin{enumerate}
  \item \textbf{Audit.} A third-party security audit of the \emph{deployed bytecode} --- not of a
  source tree that was later modified --- published in full, with every critical and high finding
  either resolved in a deployed upgrade or publicly accepted with reasoning.
  \item \textbf{Duration.} At least \num{1051200} L1 blocks ($\approx$ 12 months at the
  \SI{30}{\second} target) of continuous mainnet operation \emph{after} the audit's publication.
  \item \textbf{Exercise.} Cumulative bridged volume over that period of at least
  \num{10000000}\flux{} --- one day of a single chain's rate limit, and enough that the record reflects
  use rather than idleness.
  \item \textbf{Clean record.} Within that period: zero emergency pauses invoked, and zero upgrades
  executed to correct a contract defect. Upgrades for unrelated feature work do not reset the clock;
  an upgrade fixing a correctness bug does, because the record being demonstrated is precisely that
  no such bug remains.
  \item \textbf{Reorganization evidence.} A published log of reorganizations observed on each
  connected chain during the period, with the deepest handled on each, demonstrating that the
  finality parameters of \Cref{tab:22-recommended} were exercised rather than merely configured.
  \item \textbf{Renunciation on chain.} The upgrade path is then removed by transaction --- the
  admin role renounced, not merely unused --- so that the change is verifiable by reading the
  contract rather than by trusting a statement about it.
\end{enumerate}

Clause~5 is the one that speaks directly to the stated concern, and it is also the one most likely to
be unsatisfiable on schedule: reorganizations on the connected chains are not events the project can
schedule, and a quiet year produces no evidence. If the record is thin when the other clauses are
met, the honest options are to wait or to state that the reorganization evidence is absent --- not to
treat silence as proof.

\proposednote{The six renunciation criteria above are this paper's construction. What is settled
is their \emph{shape}: a third-party security audit plus a proven operating record across edge
scenarios including chain reorganizations \srcintent{08-answers-round-2.md, round 12}. The
quantities --- twelve months, \num{10000000}\flux{} of cumulative volume, zero defect-fixing upgrades
--- are proposed so that the stated shape becomes checkable by a third party rather than asserted;
they can be changed without changing the shape. One element is a design point rather than a number and
should not be dropped when the numbers are revised: renunciation must be executed \emph{on chain}, by
removing the admin role, so that it is verifiable by reading the contract rather than by trusting a
statement about it.}

At \qty{48}{\hour} the hypothesis (U) of \Cref{prop:worst-case} is discharged for any
$T_{\mathrm{react}} < \qty{48}{\hour}$, and the value sits at the low end of the interval this
specification gave, matching common practice for governance timelocks. Two consequences follow and
neither is enforced by code, so both are stated as conditions rather than assumed.

\emph{First, the timelock bounds the attack only if the upgrade queue is watched.} A malicious
upgrade announced and unnoticed for \qty{48}{\hour} takes effect, the rate limit is deleted, and the
loss reverts to being bounded only by the cap. Monitoring the queue is therefore a security control
of the same standing as $T_{\mathrm{react}}$, not an operational nicety.

\emph{Second, \qty{48}{\hour} is sized for the bound rather than for notice.} A timelock does two
jobs: it bounds a compromise, and it gives holders time to exit before a rule change binds them. Two
days is adequate for the first and thin for the second on a bridge carrying a nine-figure balance ---
a holder who does not look over a weekend misses the window entirely. A longer delay on changes that
affect holders, with \qty{48}{\hour} retained for security fixes, would serve both; the paper notes
this once and does not press it.
\end{remark}
\end{remark}

\paragraph{$T_{\mathrm{react}}$ is now the most important operational parameter in the bridge, and
this is not obvious from the construction.} $\mathcal{L}$ is linear in $T_{\mathrm{react}}$ with
slope $|\mathcal{C}|\,\varrho^\star / T_{\mathrm{win}}$, so over the launch set every hour of
detection-and-response latency is worth $2 \times \num{10000000}\flux / 24 \approx
\num{833333}\flux$, rising to approximately \num{1250000}\flux per hour once Solana is live
\srcinferred. Three things change that slope and the design should be read as choosing among them:
the number of live chains, the per-chain limit, and $T_{\mathrm{react}}$ itself. Several things
that look as though they should do not. The cap does not: it is \SI{25}{\day} away. The quorum does
not: it is compromised by hypothesis. The audit does not: \Cref{prop:worst-case} assumes the
contract behaves as specified, and an audit is what makes that assumption true rather than what
makes the loss smaller. Reserve custody does not, because the mint side never touches $\mathsf{R}$.
What does compress $T_{\mathrm{react}}$ is monitoring cadence,
alert routing, who is on call, whether a guardian key is reachable at 03:00 local time, and how long
a human takes to believe an alarm --- none of which are contract parameters and all of which are the
operational questions this bound turns into first-order security questions. The attestation link of
\Cref{sec:bridging-attestation} exists to compress $T_{\mathrm{react}}$, and its fail-closed silence
timeout $T_{\mathrm{silence}}$ is the mechanism that puts a ceiling on it without a human in the loop.
\planned

\begin{proposition}[Dominance over the incumbent, H0, with a number on one side]
\label{prop:dominance}
Under \Cref{tab:bridge-decided} and hypotheses (P) and (U) of \Cref{prop:worst-case}, the
replacement's worst-case loss over a reaction window is $\mathcal{L}(T_{\mathrm{react}})$, which is
\num{20000000}\flux over the launch set at $T_{\mathrm{react}} = \SI{24}{\hour}$ --- a quantity fixed
by contract constants before any compromise occurs, and the same for every compromise class that can
reach the mint path, since
Algorithm~\ref{alg:mint} applies the bucket regardless of which signatures it saw. The other
directions are bounded separately and also in advance: the release side by
$\mathsf{R}_{\mathrm{hot}}$, and the guardian and relayer roles by nothing they can move. Every
bound-consuming event is a public
transaction that Algorithm~\ref{alg:verify} reads, so H0(ii) is satisfied with $T_{\mathrm{react}}$
as the stated observability bound. The incumbent's corresponding worst case is the entire balance
held under 32 single-signature hot keys across eleven chains at the moment of compromise
\srccode{fusion/\allowbreak{}config/\allowbreak{}default.js}{342--482}: a quantity fixed by nobody, published nowhere, and
detectable only by an operator reading its own records, since an unauthorised outgoing transfer from
a hot wallet is indistinguishable on-chain from an authorised one. \planned
\end{proposition}

\begin{proof}[Proof sketch]
The replacement half is \Cref{prop:worst-case} for the mint direction, the hot/cold custody split for
the release direction, and \Cref{tab:bridge-roles} for the remaining roles, none of which can mint,
transfer or freeze. The observability half follows because each bound-consuming event is a
transaction on a public chain and Algorithm~\ref{alg:verify} reads exactly those transactions from
public state, with no team-maintained list. The incumbent half is \Cref{sec:bridging-incumbent}:
single-key in-process signing on all eleven chains, with every gate --- amount ceiling, blacklist,
rate limit --- enforced in the same process that holds the keys
\srccode{fusion/\allowbreak{}src/\allowbreak{}services/\allowbreak{}swapService.js}{150--157} \srccode{fusion/\allowbreak{}src/\allowbreak{}lib/\allowbreak{}server.js}{29,38}.
\end{proof}

\noindent\textbf{What the comparison is, stated exactly.} H0(ii) is met outright: a stated
$T_{\mathrm{react}}$ (\SI{24}{\hour}, \Cref{tab:22-recommended}) against no detection bound at all. H0(i) is met in the form that matters --- a
finite bound fixed in advance against an exposure bounded by nothing in the design --- and this is
the substantive change from earlier drafts, which could state only the shape of the bound and not its
value. It is \emph{not} a comparison of two measured numbers, because the incumbent's side has never
been measured.
\notestablished{the incumbent's current hot-wallet balances across the 32 keys are not in
the evidence corpus. \Cref{prop:dominance} therefore compares a stated finite bound against an
unbounded exposure. Should that balance ever be measured and found below
$\mathcal{L}(T_{\mathrm{react}})$, H0(i) would need restating in per-window rather than absolute
terms; the measurement is a prerequisite task, not a paper gap.}

\subsection{The attestation link}
\label{sec:bridging-attestation}

\begin{definition}[Attestation]
\label{def:attestation}
An attestation is the tuple
\begin{multline*}
\mathsf{A} = \big(\hgt,\ \mathrm{hash}_{L1}(\hgt),\ \mathsf{R}_{\mathrm{hot}},\ \mathsf{R}_{\mathrm{cold}}, \\
\{(c, h_c, \mathrm{hash}_c(h_c), M_c, \mathsf{cumMint}_c, \mathsf{cumMig}_c, \mathsf{cumBurn}_c, B_c)\}_{c \in \mathcal{C}}, \\
V_{\mathrm{lock}},\ V_{\mathrm{rel}},\ \varsigma,\ \mathsf{flags}\big)
\end{multline*}
signed by attester $j$ as $\mathrm{Sign}_{sk_j}\!\big(H(\texttt{"FLXATT1"} \| \mathsf{A})\big)$.
Every field is an objective chain fact at a stated height --- the roadmap's ``states they observed''
category \srcroadmap{flux-roadmap-public.md:135--136} --- so a false attestation is provably false to
anyone holding the relevant block. $\mathsf{flags}$ carries the signer's verdicts on the two policy
invariants the L1 cannot enforce: that every \texttt{FLXR1} range is disjoint and below $B_c$, and
that every release corresponds to a burn at depth $\ge D_c$. An attestation is \emph{valid} when
$\ge k_{\mathrm{att}}$ distinct bonded attesters sign the same $\mathsf{A}$. \planned
\end{definition}

\paragraph{Where it lives, and the cost finding that decides it.} Attestations would be published
off-chain on a public feed, with only the \emph{consequence} on-chain. Posting every attestation to
Ethereum is not viable: at one attestation per L1 block, \num{2880} per day, and
$k_{\mathrm{att}} = 9$ signatures each, the cost is approximately \num{346e6} gas per day ---
between eight and twelve Ethereum blocks' worth of gas daily for one bridge's bookkeeping
\srcinferred. At $k_{\mathrm{att}} = 5$ it is \num{276e6} and at $15$ it is \num{449e6}. The
on-chain component is therefore a single guardian contract, \texttt{Attester\allowbreak{}Guardian}, one per
chain, added to $\mathcal{G}$ alongside the four named guardians of \Cref{tab:bridge-decided} once
the attestation network exists --- which is exactly the D4 transfer path, and the only mechanism in
the design by which pause authority reaches a party outside the four: it stores the attester key set
and $k_{\mathrm{att}}$, and
exposes \texttt{pause\allowbreak{}On\allowbreak{}Attestation(A, sigs[])}, which pauses mint if and only if the signatures are
valid and $\mathsf{A}$ asserts $\varsigma < 0$ or a false policy flag. It would additionally
\emph{fail closed on silence}: if no valid attestation is submitted within $T_{\mathrm{silence}}$,
the guardian pauses mint for $\Delta_g$. That direction is deliberate. The deployed \SAS watchdog
fails \emph{open} on hash-service outage --- if both hash endpoints are unreachable it treats the
node as secure and does nothing \srccode{fluxsas/\allowbreak{}src/\allowbreak{}watchdog.h}{566--596} --- and under the team
constraint an attester outage should stop issuance rather than continue it. Burn is unaffected in
either case, so the degraded bridge is exit-only, which is the correct degradation.
\settlednote{$T_{\mathrm{att}} = \num{120}$ blocks, $T_{\mathrm{silence}} = \SI{6}{\hour}$, fail-closed adopted (\Cref{tab:22-recommended}). $T_{\mathrm{att}}$ (natural domain: one attestation per L1 block, \num{2880} per day, down
to one per hour), $T_{\mathrm{silence}}$, and whether fail-closed is adopted. Trade-off: a shorter
cadence and a shorter silence timeout detect faster and make an availability failure of the feed
into a liveness failure of the bridge.}
\settlednote{The independent mirror is mandatory (\Cref{tab:22-recommended}). The feed itself and whether a Flux-independent mirror is mandatory. IPFS is the precedent
the project already uses for parameter distribution
\srccode{fluxd/\allowbreak{}zcutil/\allowbreak{}fetch-params.sh}{17}. Trade-off: a single Flux-operated feed can withhold
but not forge; an independent mirror removes the withholding vector at operational cost.}

\paragraph{Detecting a false attestation, and what the bond is at risk for.} A false attestation is
one whose fields do not match the canonical chains at the stated heights, or whose flags assert a
policy invariant the ledgers contradict. Detection is by any party running Algorithm~\ref{alg:verify}
--- the same computation, no privileged access --- and the proof of falsity is the relevant block
header plus a Merkle inclusion or storage proof for the contradicting value. Equivocation, meaning
two different $\mathsf{A}$ for the same heights, is proven by the two signatures alone and needs no
chain data. The bond $\beta_{\mathrm{att}}$ would be at risk for signing a false $\mathsf{A}$, for
equivocation, and (minor, to keep the set honest) for sustained non-participation. The economic
condition that makes collusion unprofitable is
\begin{equation}
\label{eq:bond-condition}
k_{\mathrm{att}} \cdot \beta_{\mathrm{att}} \;>\; \varrho_c \cdot
\left\lceil \frac{T_{\mathrm{react}}}{T_{\mathrm{win}}} \right\rceil
\qquad \text{for every } c \in \mathcal{C},
\end{equation}
i.e.\ the forfeited bonds must exceed what a false ``all is well'' can cover for, which is exactly
the bound of \Cref{prop:bounded-mint}.

\medskip
\noindent\textbf{The decided rate limit makes \eqref{eq:bond-condition} very expensive, and this is
worth stating rather than discovering later.} The condition is per chain, so the per-chain limit
decision leaves it unchanged in form and in magnitude: substituting
$\varrho_c = \varrho^\star = \num{10000000}\flux$ and
$T_{\mathrm{react}} = T_{\mathrm{win}} = \SI{24}{\hour}$ gives
$k_{\mathrm{att}} \cdot \beta_{\mathrm{att}} > \num{10000000}\flux$. Using the deployed Stratus tier
collateral of \num{40000}\flux purely as a yardstick and not as a recommendation, that is
$k_{\mathrm{att}} > 250$ bonded attesters --- against $k_{\mathrm{att}} = 9$, which would put only
\num{360000}\flux at risk, short of the condition by a factor of about 28 \srcinferred. Three
responses exist and the design does not yet choose between them: raise $\beta_{\mathrm{att}}$, which
concentrates the attester set and is the failure the network exists to avoid; lower $\varrho^\star$,
which lowers the bound of \Cref{prop:worst-case} and the honest throughput ceiling together; or
accept that
the attester layer is a \emph{detector} that compresses $T_{\mathrm{react}}$ rather than a
\emph{deterrent} that makes collusion unprofitable, and say so. The third is the honest description
of any attester set small enough to operate, and it is consistent with the fail-closed pause being
the attesters' only on-chain power. \planned
\notestablished{$(n_{\mathrm{att}}, k_{\mathrm{att}})$ --- the size and threshold of the attester
set. Round 23's ruling that nothing is ever slashed changes the character of this gap rather than
merely deferring it. $\beta_{\mathrm{att}}$ is \emph{void}: a bond that cannot be forfeited is a
deposit, so there is no bond to size, and \eqref{eq:bond-condition} is withdrawn rather than left
unparameterised. What remains is an ordinary quorum choice over an accountability-based attester set
--- how many named attesters, and how many must agree --- undetermined because nobody has picked the
numbers, not because the construction is unsolved. That is a materially smaller gap than this marker
previously recorded, and it closes by decision whenever the author makes one.}

\noindent The cap, the rate limit and the window are no longer open:
$\bcap{c} = \num{250000000}\flux$, $\varrho_c = \num{10000000}\flux$ per chain and
$T_{\mathrm{win}} = \SI{24}{\hour}$ (\Cref{tab:bridge-decided}). Two constraints that the open form of
those parameters carried should be checked against the decided values rather than dropped. The first,
$\bcap{c} \ge \Lambda_c$, holds with room: the measured legacy liability on the largest chain is
\num{156217342.78}\flux on Ethereum \srcmeasured{scripts/\allowbreak{}13-issuer-balances.sh}{2026-09-03}, well
below \num{250000000}\flux. The second, $\varrho_c \ll \bcap{c}$, holds at exactly a factor of 25 and
is what makes \Cref{prop:worst-case}'s cap term inactive.

\begin{remark}[The \num{560000000}\flux{} cap is global, and no contract can enforce it]
\label{rem:global-cap}
This paper read \num{560000000}\flux{} as a per-chain contract constant $\bcap{c}$, because that is
the only reading a per-chain construction can enforce by itself. That reading is wrong: the figure is
a \textbf{global cap across all chains} \srcintent{08-answers-round-2.md, round 13} (\planned).
The correction matters, because it relocates the enforcement rather than the number.

\textbf{No contract can enforce a global cap.} A contract on chain $c$ can read its own
\texttt{totalSupply} and nothing else; it cannot observe supply on the other chains or on the L1, so
it cannot evaluate the global constraint it would need to enforce. Setting $\bcap{c} =
\num{560000000}$ on every chain yields an aggregate contract-enforced ceiling of $|\mathcal{C}|
\times \num{560000000}\flux$, which for the surviving pair is \num{1120000000}\flux{} --- twice
the intended global cap. As a backstop the constant is therefore \emph{inert}: it is set far above
any level at which it could bind.

\textbf{What actually enforces the global cap is the reserve invariant.}
\Cref{thm:reserve-dominance} establishes $\mathsf{R} \ge \sum_c M_c$: minted supply on the parallel
chains is dominated by value locked on the L1, so the global total is bounded by L1 supply and cannot
exceed it however many chains exist. This is the correct and sufficient mechanism, it is already
specified, and it is exactly the mint-and-burn discipline the team describes --- ``if we mint on some
chain, another chain is burning it or locking somewhere''
\srcintent{08-answers-round-2.md, round 8}. The global cap is a property of the reserve, not of the
contracts. That holds under honest-minority operation, assumption~(ii) of the theorem; under the
quorum compromise \Cref{prop:worst-case} models, the only contract-enforced aggregate is
$\sum_c \bcap{c}$ --- \num{500000000}\flux{} at launch, \num{750000000}\flux{} with Solana --- and
$\bcap{c}$ is not lowered for that, because the guardian's cancel over queued upgrades
(\Cref{rem:upgrade-voids}) removes the path by which the aggregate was reachable
\srcintent{08-answers-round-2.md, round 18}.

\textbf{Recommendation: set $\bcap{c}$ far lower, so it is a real backstop rather than a decoration.}
A per-chain constant is genuinely useful as defence-in-depth --- it bounds the damage from a mint-path
bug or a compromised quorum on one chain, and it does so without needing cross-chain visibility. But
it only does that work if it can bind. The largest measured per-chain liability is
\num{156217342.78}\flux{} on Ethereum (\Cref{tab:issuer-liability}), so a constant of
\num{250000000}\flux{} leaves roughly \SI{60}{\percent} headroom over the largest realistic
requirement while capping single-chain damage at under half the global supply. That is this paper's
recommendation; \num{560000000} on each chain is not a cap, it is the absence of one.
\textbf{Settled at \num{250000000}\flux{}} \srcintent{08-answers-round-2.md, round 14}
(\planned): the per-chain constant binds at a level that bounds single-chain damage below half the
global supply, while leaving roughly \SI{60}{\percent} headroom over the largest measured per-chain
liability.
\end{remark}

\paragraph{The dependency, stated honestly.} An attestation network built on today's primitives
would inherit a defect that voids its central security claim, and \Cref{sec:attested-execution}
states the same finding from the other side. The \fluxsas attestation signing seed for every purpose
is derived by HKDF over salt and domain values compiled identically into every binary; none of the
inputs depends on the TPM, on Secure Boot state, or on any per-machine value
\srccode{fluxsas/\allowbreak{}src/\allowbreak{}api\_server.h}{411--481,466--469,503--506}. A signature from that key therefore
demonstrates that the signer holds a value shipped in every copy of the software. It does not
demonstrate that a particular machine observed a particular chain state. \textbf{A key that everyone
holds cannot be slashed, and a signature it produces cannot be attributed to a bond.} The
node-identity key does mix machine-derived values, but the key protecting its random component is
itself a baked-in constant \srccode{fluxsas/\allowbreak{}src/\allowbreak{}disk\_encryption.h}{200--228}. What
\Cref{sec:attested-execution} calls Requirement~0 --- per-node attester keys, generated on the node,
whose public half is committed on the L1 in the bonded-tier registration transaction so that a
signature is attributable to a specific bond outpoint --- must be met before any of
\Cref{sec:bridging-attestation} is more than a specification. The \texttt{purpose=\allowbreak{}mesh} voucher
already present in \fluxsas with no consumer \srccode{fluxsas/\allowbreak{}src/\allowbreak{}api\_server.h}{438--445} is the
natural place for such a key to land once it is per-node. Note also that the same finding applies to
the benchmark signing path, so a bonded-tier registration must not reuse it.

\begin{openproblem}[Slashing on the \Flux L1]
\label{op:slashing}
A bond held as a plain UTXO cannot be taken from its owner by consensus, and \Flux's inherited
Bitcoin script cannot express a slashing condition. The two options are (a) a consensus change
introducing a bonded-tier transaction type with a slashing rule verified by \fluxd, or (b) bonds
escrowed in a $k$-of-$n$ that adjudicates --- which reintroduces exactly the trusted quorum the
network exists to replace. This is the load-bearing hole in the attestation link. Until it is
filled, the ``bonded'' tier is a \emph{registered} tier: attributable but not slashable, and its
attestations should be weighted in $\mathcal{G}$ accordingly.
\decided{there is no slashing, here or anywhere in \Flux{}: no mechanism destroys or forfeits a
holder's \flux{} \srcintent{08-answers-round-2.md, round 23}. That closes the construction question
by removing the mechanism rather than by supplying one, and the consequence has to be carried rather
than absorbed. \textbf{An attester bond that cannot be forfeited is a deposit, not a stake}, so
\eqref{eq:bond-condition} cannot rest on it and the cryptoeconomic form of the security argument is
unavailable. What remains is accountability: a named, non-anonymous attester set whose members can be
excluded from future participation and who answer for their attestations personally, in the same
shape as the key custody of \Cref{sec:model}. That is a materially weaker claim than a bonded one,
and this paper does not dress it as equivalent --- an attester who lies loses future income and
standing, not principal, and what bounds a dishonest quorum is the rate limit and the reserve
invariant of \Cref{con:bridge}, not a forfeiture.}
\end{openproblem}

\subsection{Why not the alternatives}
\label{sec:bridging-alternatives}

\paragraph{Atomic swaps --- the question is closed, not open.} An HTLC atomic swap exchanges an
asset one party holds on chain A for an asset a counterparty holds on chain B, atomically, with no
custodian \citep{herlihy2018atomic}. \Flux's inherited script carries the primitives. But an atomic
swap cannot do what a bridge does: \textbf{it cannot issue FLUX on Ethereum}, only exchange FLUX for
something that already exists there. Everything else is secondary and still disqualifying: it
requires a counterparty holding the other asset, a discovery layer to find them, both parties online
for the duration, and it hands the initiator a free option to abort if the price moves during the
timelock. The correct conclusion is not that atomic swaps are a rejected bridge design but that they
address a different problem --- the \emph{exchange} of FLUX against ETH, USDC or another asset
against a market maker, which is \FusionX's product (\Cref{sec:bridging-incumbent}), not the
bridge's. Atomic swaps are a settled construction in the literature \citep{herlihy2018atomic}; the boundary drawn here is a product decision, not a doubt about the primitive.

\paragraph{Light-client verification --- not constructible for \Flux, and here is exactly why.} The
trust-minimizing answer would be an Ethereum contract that verifies \Flux headers directly, as IBC
does between Tendermint chains \citep{cosmosibc} and as the interoperability literature identifies as
the strongest available class \citep{zamyatin2021sok}. It is not constructible from the current
\PoN header, for three compounding reasons:

\begin{enumerate}
\item The \PoN block header carries \texttt{nodes\allowbreak{}Collateral} and \texttt{vchBlockSig} but
      \textbf{no commitment to the confirmed node set}
      \srccode{fluxd/\allowbreak{}src/\allowbreak{}primitives/\allowbreak{}block.h}{44--45,64--96}, so a verifier cannot check that a
      block's signer was eligible without the full \FluxNode registry --- which is to say, without
      being a full node.
\item Chainwork is a fixed $2^{40}$ per block regardless of difficulty target
      \srccode{fluxd/\allowbreak{}src/\allowbreak{}pow.cpp}{284--290}, so fork choice degenerates to a block-count race and
      there is no cumulative-work heuristic by which a light client could prefer one chain over
      another.
\item The 2-of-4 emergency-block path can produce a valid block outside the normal rules, with no
      time or stall-detection gating \srccode{fluxd/\allowbreak{}src/\allowbreak{}emergencyblock.cpp}{183--191}, so no
      light-client rule derived from \PoN eligibility is sound; a perfect light client would still
      be trusting those four keys.
\end{enumerate}

\noindent A future header commitment to the deterministic node list's Merkle root would reopen the
question --- after which verifying \num{2880} headers a day on Ethereum becomes a cost question, one
\texttt{ecrecover} plus a $\lceil \log_2 6489 \rceil = 13$-hash Merkle proof per header
\srcinferred --- and \Cref{sec:chain-evolution} should treat it as a candidate for the same network
upgrade that carries collateral maturity. \textbf{This is the paper's clearest instance of vertical
integration cutting both ways.} Controlling the L1 is what lets \Flux put an entitlement memo in an
\texttt{OP\_RETURN}, freeze a halving by hand, and change the reorg limit to launch \PoN; it is also
what produced a header format that forecloses the most trust-minimizing bridge architecture
available to other chains. Owning the stack does not automatically mean the stack was designed for
what you later want from it. The same tension runs through \Cref{sec:fluxos}. The reverse direction
--- \fluxd verifying Ethereum --- would need an EVM and BLS light client inside the daemon, a
consensus change of a different order.

\paragraph{Threshold-signature and MPC bridges.} One aggregate key, signed by a threshold of parties
holding shares. \emph{Trust assumption:} fewer than the threshold of share-holders collude, and the
share distribution is what it is claimed to be. \emph{Failure mode:} on-chain the contract sees a
single signer, so who participated is invisible and the claimed distribution is unfalsifiable from
outside; if the shares sit on one organisation's servers the threshold is cosmetic, which is
precisely the Multichain outcome (\Cref{sec:bridging-incidents}). \emph{Why not chosen:}
attributability is a hard requirement here (H3), and the construction additionally declines to rely
on the threshold alone --- the contract bounds the damage whatever the signing scheme. A TSS could
still be used \emph{inside} one signer organisation to protect its own key.

\paragraph{External validator sets.} A named set signs mint messages; the contract mints on
$k$ signatures. \emph{Trust assumption:} fewer than $k$ collude or are compromised.
\emph{Failure mode:} the contract does whatever $k$ signatures say, without bound --- Wormhole's 19
guardians at 13 and Ronin's 9 validators at 5 both failed this way, for different reasons.
\emph{Why not chosen, and the honest statement:} at the signing layer this construction \emph{is} a
validator-set bridge with a smaller set than either --- 4 signers at a threshold of 3, against
Wormhole's 19 at 13 and Ronin's 9 at 5 --- and pretending otherwise would be dishonest. What differs
is not the size or the independence of the set but that the cap, the per-chain rate limit, the
1-of-4 guardian pause and the admin timelock are intended to make $k$-collusion bounded and slow
rather than total and instant: \num{20000000}\flux per \SI{24}{\hour} of reaction time over the
launch set (\Cref{prop:worst-case}) instead of the pooled value. That sentence is true only under
the timelock condition of \Cref{rem:upgrade-voids}; an upgradable contract without one puts this
construction back in the same row as the systems it is being distinguished from, which is why that
item was the section's load-bearing open one until \Cref{rem:timelock-decided}. The
collateral-backed alternative in the literature --- overcollateralized vaults with automatic
liquidation \citep{zamyatin2019xclaim} --- removes the trusted set at the cost of locking capital
worth more than the bridged value on the destination chain, and of a price oracle for the
collateral-to-asset ratio. That trade is available and was not taken; the reason is that it
substitutes an oracle dependency and a capital cost for a signer dependency, and does not remove the
need for the L1-side release quorum, which is where \Cref{sec:bridging-open} shows the residual
trust actually sits.

\paragraph{Canonical rollup bridges and IBC, for completeness.} Optimism's and Arbitrum's canonical
bridges derive security from the L1 verifying the L2 through fraud or validity proofs with a
challenge window; neither \Flux nor Ethereum can verify the other, so this is not available.
Their \emph{operational} pattern --- guardian pause, security-council timelock, no silent upgrades
--- is adopted wholesale. IBC requires cheap, fast-final header verification on both ends
\citep{cosmosibc}; \Flux has neither cheap header verification nor a node-set commitment, so it is
not applicable without the header change above.

\paragraph{Issuer-attested burn-and-mint, which is the model chosen.} Circle's CCTP signs burn
messages with an issuer-operated attestation service; destination contracts mint only against
attested, nonce-bound messages, and no pooled liquidity exists to steal \citep{cctp}. Its trust
assumption is the issuer's attester key --- which, for USDC, belongs to the same party that can mint
USDC anyway, so the bridge adds no new trust. \Flux is the closest analogue and this is the chosen
model. Two differences follow from the fact that \Flux's attesters are \emph{not} already trusted for
the whole supply: an explicit $k$-of-$n$ with on-chain attribution rather than a single issuer
signature, and contract-level cap and rate limit, which CCTP does not need.

\subsection{The launch trust assumption: one quorum, bounded by a rate limit}
\label{sec:bridging-assumption}

\begin{assumption}[Single attributable quorum, bounded by contract]
\label{as:admin}
At launch the MINT, RELEASE and ADMIN authorities of \Cref{tab:bridge-roles} are held by \emph{one}
set: $n_M = n_R = n_A = 4$ and $k_M = k_R = k_A = 3$, the four keys held by the four named
co-founders --- Tadeas Kmenta, Daniel Keller, Jeremy Anderson, Valter Silva
\srcintent{08-answers-round-2.md, round 5} --- with $\mathcal{G}$ the same four individuals but at a
threshold of \textbf{one} (\Cref{tab:bridge-decided}). The assumption is that fewer than three of
those four keys are compromised or collude. It is one assumption, not four independent ones. What
supports it:
(i) each key is per-signer, so every authorisation is attributable on-chain to the key that produced
it, and the signer set is published;
(ii) each key is generated and held on dedicated signing hardware and never on a host that runs a
relayer, an RPC, or the incumbent service --- co-locating the key with the process that decides is
the incumbent's defining mistake (\Cref{sec:bridging-incumbent});
(iii) $k/n = 3/4 \ge \lceil 2n/3 \rceil / n$, so the threshold ratio is at or above the shape the
requirements recommended.
What is \emph{not} satisfied, and was written into the earlier form of this assumption precisely
because it matters: signers are named individuals rather than independent named organisations; one
organisation's founders hold every key in every set, so no clause bounding a single organisation's
share of a quorum can hold; the four memberships coincide exactly; and jurisdictional spread is
whatever those four individuals happen to have rather than a design parameter.
If the assumption fails, the contract still bounds unbacked issuance to
$\mathcal{L}(T_{\mathrm{react}})$ (\Cref{prop:worst-case}) --- the pause remains reachable at 1-of-4
(\Cref{cor:pause}), and the bound additionally requires the upgrade timelock of
\Cref{rem:upgrade-voids} --- and the L1 side to $\mathsf{R}_{\mathrm{hot}}$. \planned
\end{assumption}

\begin{remark}[The separation of powers is not available at launch, and saying otherwise would be false]
\label{rem:sets-coincide}
Three compromised keys can, in one sequence and with no second party: authorise mints up to the rate
limit; co-sign L1 spends of $\mathsf{R}_{\mathrm{hot}}$; queue an admin change raising $\bcap{c}$ and
$\varrho_c$ and replacing $Q_M$ with keys of their own; and \emph{upgrade the contract}, which
reaches every rule in \Cref{tab:bridge-roles} at once including the guardian rule
(\Cref{rem:upgrade-voids}). What they cannot do is throw the pause on the honest guardian's behalf,
prevent that guardian from throwing it, or prevent that guardian from cancelling what they queue,
because the pause and the cancel are 1-of-4 (\Cref{cor:pause}) --- and
that single exception is what the whole quantitative claim of this section rests on. The $\Delta_A$
notice period survives for parameter changes --- an admin change is still visible for $\Delta_A$
before it executes --- but the exit it is supposed to protect runs through a release quorum
consisting of the same three keys, so a holder who burns during the notice period has a public claim
and the same adversary decides whether it is paid (\Cref{op:liveness}). The total exposure of a
3-key compromise is therefore $\mathcal{L}(T_{\mathrm{react}}) + \mathsf{R}_{\mathrm{hot}}$, with
the cold reserve as the last line and $\sum_c \bcap{c}$ as the ceiling if the upgrade path is not
timelocked or a queued upgrade is not cancelled within $\Delta_A$. This is stated as a property of
the launch configuration, not as a hypothetical: the
earlier draft of this remark described the collapse as the reason separation was a requirement, and
the configuration that has been decided is the collapsed one everywhere except the pause.
\end{remark}

\paragraph{Why this is defensible, on the team's own constraint --- and what would make it not.} The
governing constraint of this section is that \emph{a decentralized design that is exploitable is
worse than the custodial incumbent} \srcintent{03-parallel-assets.md}. Taken seriously, it settles
the launch question. A more distributed signer set would have to be assembled from parties who can be
held to their signatures; the mechanism for holding them --- a bond that can be forfeited --- does
not exist on this chain (\Cref{op:slashing}), and the attestation keys that would identify them are
today fleet-shared rather than per-node (\Cref{sec:bridging-attestation}). A fifteen-key set
recruited without either property is not obviously more decentralized than four named founders: it is
the same trust assumption spread across more people who are harder to identify and cannot be
penalised, which is the shape the Ronin and Multichain rows describe
(\Cref{tab:bridge-incidents}). Against that, the decided
configuration takes the trust concentration as given and spends its design effort on making the
consequence bounded and public: the contract refuses beyond $\varrho_c$ and $\bcap{c}$ whatever the
three keys say, and anyone can run Algorithm~\ref{alg:verify} without asking the signers for
anything. \textbf{The accurate description is a custodial bridge with contract-enforced bounds and
public verification, moving toward autonomy} --- not a decentralized bridge, and not a separation of
powers. The intent on record is explicit about the direction and about the fallback: full autonomy
afterwards, and ``if the autonomous design is not ready in time, ship 3-of-4 and take a second fork
for the autonomous bridge'' \srcintent{08-answers-round-2.md}. Two things would make the position
indefensible rather than defensible, and both are checkable. \textbf{If the upgrade path ships
without a timelock satisfying $\Delta_A > T_{\mathrm{react}}$, the bound is the cap rather than
\num{20000000}\flux, and the argument above does not survive the substitution}
(\Cref{rem:upgrade-voids}) --- the contract that ``refuses beyond $\varrho_c$ and $\bcap{c}$
whatever the three keys say'' is then a contract those same three keys can replace before it
refuses anything. And if the second fork does not follow, what is described here as a launch
configuration becomes the design, at which point the honest comparison is no longer against the
incumbent but against a competently run custodian. \planned

\paragraph{The path off the single quorum, specified as requirements because the construction is
open.} The intent is autonomy, and no design for it exists in the evidence corpus. What replaces the
3-of-4 must satisfy, at minimum: \textbf{(R1)} authorisation attributable to a set whose membership
is derivable from public chain state rather than from a published list, so that an observer can
recompute who was entitled to sign at a height; \textbf{(R2)} a penalty that can actually be applied
to a misbehaving member, which on this L1 means \Cref{op:slashing} is answered first; \textbf{(R3)}
per-member keys generated on the member's own machine, which means \Cref{sec:attested-execution}'s
Requirement~0 is met first; \textbf{(R4)} liveness under member churn, since a set derived from the
node registry changes between blocks while an L1 P2SH redeem script does not; and \textbf{(R5)} a
migration path from the 3-of-4 that does not require holders to act, which under D3 means the ADMIN
role rotating $Q_M$ and $Q_R$ rather than a contract replacement. The interfaces are already fixed by
\Cref{con:bridge} and do not change: the MintAuth digest, the $k$-of-$n$ check inside
Algorithm~\ref{alg:mint}, and the P2SH release path. \textbf{The construction that meets R1--R5 is
open}, and R2 and R4 are the binding pair on this chain specifically: R2 requires a consensus change
(\Cref{op:slashing}), and R4 is in direct tension with a release path whose signer set is frozen into
a P2SH redeem script at funding time and capped at 15 keys.
\proposednote{The autonomous replacement for the 3-of-4 quorum --- requirements R1--R5 above,
interfaces unchanged --- is not constructed here, and the reason is now specific rather than
general: it is gated on the same slashing problem that leaves $\beta_{\mathrm{att}}$ undeterminable,
because any replacement that removes the quorum must put something forfeitable in its place. What
\emph{is} settled is the intent and the exit condition: the quorum is an explicitly temporary
bring-up mechanism, and immutability is the stated destination, reachable through the criteria above
\srcintent{08-answers-round-2.md, rounds 11--12}.}

\subsection{Migration and consolidation}
\label{sec:bridging-migration}

\Cref{sec:parallel-assets} inventories ten chains carrying FLUX today. All ten legacy contracts
would be retired, including the legacy Ethereum, BSC and Solana tokens, because the survivors are
\emph{relaunched} rather than kept \srcintent{03-parallel-assets.md}; three new deployments would
follow, ETH and BSC at launch and Solana later. \planned

\begin{remark}[Two migration paths, and ``surviving'' is not the same set as ``carried'']
\label{rem:two-paths}
The distinction matters more than the phrasing suggests, and conflating the two sets is the easiest
mistake a reader could make here. \textbf{Only Ethereum and BSC are snapshot-and-distributed}: their
balances are captured at $\hsnap$ and reissued $1{:}1$ into the new contracts, coordinated
with exchange partners and ecosystem integrators, and their holders take no action at all
\srcintent{evidence/design/08-answers-round-2.md}. \textbf{Every other chain redeems through the
incumbent bridge} to the main chain during the six-month window, with unclaimed balances passing to
the Foundation at close.

\textbf{Solana sits in both descriptions and belongs to the second.} It is a surviving chain in the
sense that a new Solana deployment follows the launch pair --- but its \emph{legacy} holders are not
carried across. They redeem to the main chain like the holders of any retired chain, and re-enter the
new Solana deployment later if they choose. So the snapshot-carried set is exactly
$\{\mathrm{ETH}, \mathrm{BSC}\}$, and the redeeming set is the other eight chains, Solana included.

This has a measurement consequence \Cref{sec:bridging-migration} states below rather than leaves
implicit. An earlier draft of this remark cautioned that Solana might prove the largest single
component of the redeeming set, on the reasoning that its legacy \texttt{totalSupply} measures
\num{59999990.375} \srcmeasured{api.mainnet-beta.solana.com \texttt{getTokenSupply}}{2026-09-03}.
That caution was wrong, and measuring it is what showed so: \SI{99.4}{\percent} of that supply is
issuer-held --- \num{54000000} in the \texttt{LOCKED} account alone --- leaving
\num{373067.20}\flux{} outstanding, the \emph{smallest} component of the redeeming set rather than
the largest \srcmeasured{scripts/\allowbreak{}15-nonevm-liability.sh}{2026-09-04}. Supply is not liability, and
the gap between them on Solana is a factor of \num{161}. The general point survives the specific
error: an unmeasured chain contributes an unknown, not a small, quantity, and the direction of the
surprise is not predictable in advance --- which is why the redesign requires
$\Lambda_c$ measured rather than assumed.
\end{remark}

\paragraph{The measurement prerequisite: partly done.} Before any migration mint, $\Lambda_c$ must be
measured for each of the ten legacy assets at a published snapshot height $\hsnap$ by
transfer-log replay, and the available backing $R_0$ (\Cref{def:reserve}) must be established and
shown to satisfy $R_0 \ge \Lambda$. $\Lambda_c$ is now measured for all ten legacy assets
(\Cref{tab:issuer-liability}), totalling \num{223027781.21}\flux outstanding \shipped
\srcmeasured{scripts/\allowbreak{}13-issuer-balances.sh}{2026-09-03};
\srcmeasured{scripts/\allowbreak{}15-nonevm-liability.sh, scripts/\allowbreak{}16-remaining-liability.sh}{2026-09-04}. The
distribution matters more than the total: \num{216800880.41}\flux --- \SI{97.21}{\percent} --- is on
the two surviving chains and is carried by snapshot, and the eight retiring chains that require
holder action carry \num{6226900.80}\flux between them --- \SI{2.79}{\percent} of the total.
Kadena required a different operand: its Pact module exposes no total-supply accessor, so supply was
recovered by summing the \texttt{ledger} table across all twenty Chainweb chains
\srcmeasured{scripts/\allowbreak{}19-kadena-ledger.js}{2026-09-04}. What remains needed before a migration mint is
confirmation that reserves cover \num{223027781.21}\flux. H8 forbids opening migration at 1:1 on any
chain until both the measurement and its coverage confirmation are in hand for that chain.
\notestablished{the composition and balance of $R_0$ against the now fully-measured
\num{223027781.21}\flux of legacy liability; and per-chain holder counts are not established anywhere in the evidence
corpus. This is a prerequisite task blocking \Cref{sec:bridging-migration} and
\Cref{sec:parallel-assets}, not a gap the paper can close by argument. If the measurement shows a
shortfall, the migration cannot open at 1:1 until it is covered, and the paper must say so.}

\paragraph{Two migration mechanisms, not one, and the difference decides who is exposed.} The
decided plan uses \emph{different} mechanisms for the surviving and the retiring chains
\srcintent{08-answers-round-2.md, round 7}, and conflating them is what produced the
two-orders-of-magnitude error in earlier drafts of this subsection and of
\Cref{sec:parallel-assets}. \textbf{Surviving chains are carried by snapshot; retiring chains are
redeemed by the holder.} The first requires no holder action and therefore no holder is at risk of
the forfeiture term; the second requires an act within a six-month window and therefore every
non-acting holder is. \planned

\smallskip
\noindent\emph{Snapshot-and-carry (ETH, BSC --- surviving).} Because the surviving contracts are
\emph{relaunched} rather than kept (\Cref{sec:parallel-assets}), holders of the legacy token are
moved to the new contract by the issuer rather than by themselves. The holder set is reconstructed
off-chain at a published height $\hsnap$ by \texttt{Transfer}-log replay --- the same
computation the team's own crawler already performs on four of the five EVM chains
\srccode{holder-snapshot/\allowbreak{}eth.js}{63--88} --- and written into the new contract by
Algorithm~\ref{alg:carry}, which consumes the same $\Lambda_c$ migration allowance that
\texttt{mintMigrated} does, so bracket~(1) of \Cref{thm:reserve-dominance} and the conservation
argument are unchanged. What changes is only who acts. The exercise is coordinated in advance with
exchange partners and ecosystem integrators \srcintent{08-answers-round-2.md, round 7}, which is not
a courtesy: a custodian's users hold no address of their own on the legacy contract, so a credit to
the custodian's address is the only carry available for them and the custodian's internal ledger is
what makes it reach a user.

\begin{algorithm}[H]
\caption{Snapshot carry onto the relaunched contract for a surviving chain $c$. \planned}
\label{alg:carry}
\KwIn{The published snapshot height $\hsnap$; the published issuer-held address set
$\mathcal{I}_c \subseteq \mathcal{A}_c$, whose balances are excluded from $\Lambda_c$
(\Cref{tab:issuer-liability}) \srcintent{08-answers-round-2.md, round 18}; the balance map
$\mathsf{bal}: \mathcal{A}_c \to \mathbb{Z}_{>0}$ reconstructed by \texttt{Transfer}-log replay of
the legacy contract at $\hsnap$, over the address set $\mathcal{A}_c$; the migration
allowance $\Lambda_c$; the new contract with $M_c = 0$, public mint paused.}
\KwOut{Each $a \in \mathcal{A}_c \setminus \mathcal{I}_c$ credited exactly $\mathsf{bal}(a)$, then
$\mathsf{snapshotComplete}_c \leftarrow \mathsf{true}$.}
\KwBrInv{$\mathsf{cumMig}_c = \sum_{a \,\mathrm{credited}} \mathsf{bal}(a) \le \Lambda_c$ at every
point, and no $a$ is credited twice --- enforced by a per-address
$\mathsf{carried}_c \subseteq \mathcal{A}_c$ set in contract storage, not by the caller's
bookkeeping. Transfers and public mint stay disabled until
$\mathsf{snapshotComplete}_c$, so no partially-carried supply is ever tradeable.}
\KwBrFail{A batch reverts atomically and is resubmitted; a batch naming an already-carried address
reverts rather than double-crediting. If the replay is wrong, the error is public and diffable
against the legacy contract at $\hsnap$ by anyone, which is the only reason this path is
acceptable at all --- it is an issuer-executed mint whose correctness is checkable by third
parties.}
\BlankLine
publish $\hsnap$, $\mathcal{I}_c$, the reconstructed $\mathsf{bal}$ and its Merkle root before executing\;
\ForEach{batch $\mathcal{K} \subseteq \mathcal{A}_c \setminus (\mathsf{carried}_c \cup \mathcal{I}_c)$}{
  \lIf{$\mathsf{cumMig}_c + \sum_{a \in \mathcal{K}} \mathsf{bal}(a) > \Lambda_c$}{\Return revert (allowance)}
  credit each $a \in \mathcal{K}$; $\mathsf{carried}_c \mathrel{\cup}= \mathcal{K}$;
  $\mathsf{cumMig}_c \mathrel{+}= \sum_{a \in \mathcal{K}} \mathsf{bal}(a)$\;
}
\lIf{$\mathsf{carried}_c \ne \mathcal{A}_c \setminus \mathcal{I}_c$}{\Return abort (incomplete carry; do not enable transfers)}
$\mathsf{snapshotComplete}_c \leftarrow \mathsf{true}$; enable transfers and public mint\;
\end{algorithm}

\begin{remark}[The snapshot carry --- two elements settled, one specified here]
\label{rem:snapshot-settled}
\textbf{(i) Contract-held balances: resolve through to the beneficial holder.} A legacy balance held
by an AMM pair, a lending market or a third-party bridge cannot be carried by crediting the
contract's address --- the new token is a different address and the position does not follow it. The
settled treatment is to read the position and credit the actual holder: ``from LP, staking we can see
from the LP positions who is the actual holder and that we will distribute to the actual holder like
we did in the past'' \srcintent{08-answers-round-2.md, round 9} (\planned). This is a stronger
commitment than any of the three options previously listed, and the team reports precedent for
having done it before. Two obligations follow and the paper states them: the resolution is
\emph{per-protocol} work --- an AMM pair, a lending market and a bridge each expose beneficial
ownership differently, and some expose it only at a block height, not retrospectively --- and it must
be published as a list before $\hsnap$, because a holder cannot check a resolution they
cannot see. \decided{the contract-held positions are not measured in this paper, and that is a scoping choice
rather than a limit. Balances held in liquidity-pool and staking contracts on Ethereum and BSC are
public on-chain, readable by the same holder-enumeration method already used for the parallel-asset
snapshots of \Cref{sec:parallel-assets}, and the team has performed that measurement before
\srcintent{08-answers-round-2.md, round 24}. What matters for the migration argument is that the
quantity is \emph{obtainable}: a holder-address snapshot alone will miss these positions, and the
carry of \Cref{sec:bridging-migration} must therefore resolve them per protocol before
$h_{\mathrm{snap}}$ --- work this paper specifies but does not perform, because running it belongs to
the migration rather than to the survey.}

\textbf{(ii) Announcement: $\hsnap$ is announced in advance} \srcintent{08-answers-round-2.md,
round 9} (\planned). This is the right choice and the paper does not hedge it: it lets holders
arrange their balances, and it is the only option compatible with publishing a contract-held
resolution list beforehand. It does admit the griefing vector previously noted --- an adversary may
fragment a balance across many addresses to inflate the carry's gas cost --- and that vector has no
defence in this construction. It is cheap for the adversary and merely expensive for the issuer, so
the paper records it as an accepted cost rather than an open problem, and \Cref{rem:batch-carry}
sizes the exposure.

\textbf{(iii) Batch size and gas: specified below} (\Cref{rem:batch-carry}), as a recommendation.
\end{remark}

\begin{remark}[Batch sizing for the carry --- recommendation]
\label{rem:batch-carry}
The carry is $|\mathcal{A}_c|$ storage writes and the issuer pays the gas. A cold \texttt{SSTORE} on
an EVM chain costs \num{20000} gas, plus per-call overhead and calldata; at a \num{30000000} gas
block limit, a batch of \num{500} credits consumes roughly \SI{35}{\percent} of one block and leaves
headroom for the base transaction and calldata \srcinferred. The paper recommends
\textbf{batches of \num{500} recipients}, submitted sequentially with each batch's completion
confirmed before the next is sent, and an idempotent \texttt{claimed} flag per recipient so a
reverted or re-sent batch cannot double-credit.

The trade-off is stated exactly: larger batches are cheaper per recipient and put more value at risk
in a single reverted transaction. At \num{500} the worst case is one reverted batch of \num{500}
credits, which is recoverable by resubmission because the flag makes the operation idempotent; the
gas is lost, the entitlements are not. Smaller batches buy little, because the per-recipient cost is
dominated by the \texttt{SSTORE}, not the per-transaction overhead.
\end{remark}

\smallskip
\noindent\emph{Holder-initiated redemption (Avalanche, Base, Polygon, and the non-EVM retirees).}
The holder must act, within the window, or lose the claim. This is the path the forfeiture term
applies to, and it is the only one it applies to. \emph{Retired EVM chains} (the disposal construction
this section first specified; it will not be deployed, \Cref{rem:no-legacysink}): a
\texttt{LegacySink} contract per chain holds legacy tokens forever with no withdraw function --- the
legacy ERC-20 has no \texttt{burn} \srccode{flux-eth-bsc/\allowbreak{}flux.sol}{6--13}, so a sink is the only
irreversible disposal --- and
\texttt{deposit(v, flux\allowbreak{}Addr)} emits an event the RELEASE quorum treats exactly as a Burn --- depth
$D_c$, contiguous nonce ranges, \texttt{FLXR1} memo --- paying L1 FLUX 1:1. \emph{Retired non-EVM
chains} (legacy Solana, Kadena, Tron, Algorand, Ergo): the holder makes the legacy balance provably
unspendable and binds the disposal to a \texttt{fluxAddr}; Solana has a self-custodial SPL
\texttt{Burn}, and Ergo admits a \texttt{sigma\allowbreak{}Prop(false)} box.
\settlednote{Moot: no disposal contract is deployed on any chain (\Cref{rem:no-legacysink}), so no sink construction is required for Kadena, Tron or Algorand. Retained as a pointer because \S\ref{sec:22:recommended} records what the three constructions would be if a sink is ever wanted.}

\paragraph{The decided migration terms.} Four terms are settled
\srcintent{08-answers-round-2.md, rounds 5 and 7}, and they replace the window and disposition
parameters this subsection previously carried as open. \planned

\begin{enumerate}
\item \textbf{Effective date.} The contract changes are prepared now and take effect at PA Depletion,
      $\hpa = \num{3466630}$ (approximately 2027-03), \emph{simultaneously} with the swap contract
      and with \PNR activation (\Cref{sec:pnr}, \Cref{sec:migration}). Nothing takes effect earlier
      and there is no phased cutover.
\item \textbf{Redemption window.} Six months from $\hpa$ --- \SI{182.5}{\day} nominal, approximately
      \num{525600} L1 blocks at the \PoN target spacing \srcinferred --- so
      $h^{\mathrm{mig}}_{\mathrm{close}} \approx \num{3992230}$. During the window, redemption
      back to main-chain FLUX runs at 1:1 through the \textbf{existing bridge}.
\item \textbf{The redemption route is one-directional.} During the window the existing bridge permits
      migrating \emph{to} the main chain and claiming \emph{on} the main chain, and nothing else:
      holders can move to the main chain and nowhere else
      \srcintent{08-answers-round-2.md, round 7}.
\item \textbf{Unclaimed balances.} At window close, \textbf{unclaimed balances pass to the Flux
      Foundation to fund development.}
\end{enumerate}

\noindent The fourth term is stated here in those words because it is a transfer from non-acting
holders to the issuer, and naming it, its window and its destination is what makes it defensible.
Unclaimed value must belong to someone; a stated destination is better than an unstated one; and the
alternative --- an open migration authority honoured forever --- keeps the largest single mint
authority in the design alive indefinitely, which is a cost paid by everyone rather than by the
non-acting holder. What passes is not the legacy tokens themselves, which sit on foreign chains and
cannot be recalled by anyone (see \emph{The long tail} below); it is the residual main-chain backing
$\Lambda - \sum_c \mathsf{cumMig}_c - V^{\mathrm{mig}}_{\mathrm{rel}}$, which is an explicit public
number at every height and is computed by Algorithm~\ref{alg:verify} without privileged data.

\medskip
\noindent\textbf{The amount actually exposed to the window, which is two orders of magnitude below
the measured liability.} The forfeiture term applies only to the holder-initiated path, so the
exposed quantity is the outstanding balance on the \emph{retiring} chains and not the ten-chain
total. Of the \num{223027781.21}\flux measured across all ten chains
\srcmeasured{scripts/\allowbreak{}13-issuer-balances.sh}{2026-09-03}:

\begin{table}[H]
\centering\footnotesize
\begin{tabular}{@{}llr@{}}
\toprule
chains & migration path & outstanding (\flux) \\
\midrule
Ethereum, BNB Chain & snapshot-and-carry; \textbf{no holder action}; not exposed & \num{216800880.41} \\
\midrule
Base & holder-initiated redemption; \textbf{exposed to the window} & \num{2528858.43} \\
Polygon & \textit{as above} & \num{647928.73} \\
Avalanche C-Chain & \textit{as above} & \num{628582.62} \\
Ergo & \textit{as above} & \num{486746.62} \\
Algorand & \textit{as above} & \num{389643.01} \\
Solana & \textit{as above} & \num{373067.20} \\
Kadena & \textit{as above} & \num{798207.66} \\
Tron & \textit{as above} & \num{373866.53} \\
\midrule
eight retiring chains & --- & \textbf{\num{6226900.80}} \\
all ten chains & --- & \num{223027781.21} \\
\bottomrule
\end{tabular}
\caption{Exposure to the six-month window and the unclaimed-balance term, by migration path.
Surviving-chain holders are carried across by \Cref{alg:carry} and take no action, so they are not
at risk of forfeiture; only the retiring chains are. Figures are
\srcmeasured{scripts/\allowbreak{}13-issuer-balances.sh}{2026-09-03} and
\srcmeasured{scripts/\allowbreak{}15-nonevm-liability.sh, scripts/\allowbreak{}16-remaining-liability.sh,
scripts/\allowbreak{}18-tron-final.sh, scripts/\allowbreak{}19-kadena-ledger.js, scripts/\allowbreak{}20-kadena-issuer.js}{2026-09-04},
the exact measured values of \Cref{tab:issuer-liability}, not rounded to make the point. All ten
legacy assets are measured; no component of this table is an estimate.}
\label{tab:migration-exposure}
\end{table}

\noindent So the amount exposed to the six-month window and the unclaimed-balance term is
\textbf{\num{6226900.80}\flux across the eight retiring chains}. That is \SI{2.79}{\percent} of
the \num{223027781.21}\flux measured across all ten chains: the overwhelming majority of parallel-asset supply is
issuer-held and never left the treasury, and the overwhelming majority of what did is on the two
chains that survive and are carried across without holder action. \textbf{The fraction of that which fails to migrate
within six months is unknown, and this paper does not estimate it}: there is no holder-count or
dormancy measurement for any of the ten legacy assets in the evidence corpus, and a number produced
without one would be a guess placed next to a real measurement.
\srcintent{08-answers-round-2.md, round 7}

\paragraph{Three consequences of routing redemption through the incumbent, one of which is a real
mitigation.} First, the incumbent's custody model is not retired at $\hpa$ but at window close six
months later. For that interval the bound of \Cref{prop:worst-case} applies to the new
lock-and-mint path and \emph{not} to the redemption path, whose exposure remains the balance under
the 32 hot keys with no bound and no on-chain attribution (\Cref{prop:dominance}). The section's
dominance claim is therefore a claim about the steady state after window close, and about the new
path from $\hpa$; it is not a claim about the window. \textbf{Second, and this materially narrows
the first: the route is one-directional} \srcintent{08-answers-round-2.md, round 7}. The incumbent's
residual role during the window is an \emph{exit to the main chain}, not general two-way traffic, so
the swap engine's chain-pair matrix, its several disagreeing confirmation tables and its
finished-before-sent state machine (\Cref{sec:bridging-incumbent}) are not all in use for it. The
exposure that remains is real --- the same single hot key per chain authorises the outbound main-chain
payment --- but it is the exposure of one direction rather than of eleven chains' worth of pairs, and
the section states it that way rather than repeating the stronger earlier framing. Third, the
\texttt{LegacySink}-plus-RELEASE-quorum path specified above becomes a \emph{disposal} mechanism for
retired units rather than the redemption route, which changes what has to exist at $\hpa$ and what
may follow later.
\begin{remark}[\texttt{LegacySink} is not deployed --- settled]
\label{rem:no-legacysink}
\texttt{LegacySink} was a construction proposed by this paper, not an existing component: a contract
with a \texttt{deposit} function and no withdraw function, supplying the irreversible disposal that
the legacy ERC-20 lacks, since it has no \texttt{burn}
\srccode{flux-eth-bsc/flux.sol}{6--13}. Its purpose was to prevent a unit being redeemed for L1 FLUX
and then sold on to a buyer who does not know the claim has already been made.

\textbf{It will not be deployed} \srcintent{08-answers-round-2.md, round 14} (\planned), and the
reasoning is a cost comparison the decision to route redemption through the incumbent makes decisive.
Building, auditing and deploying a disposal contract on eight retiring chains --- including the three
non-EVM constructions that have no common shape (\S\ref{sec:22:recommended}) --- is substantial work
whose only benefit is preventing resale of already-claimed units. Against it: the entire retiring
exposure is \num{6226900.80}\flux{} (\Cref{tab:migration-exposure}), the contracts are being
abandoned in any case, and audit A1 would have to cover a contract deployed only to be discarded.

\textbf{The mitigation is therefore the public record alone, and the paper says so rather than
implying more.} A redemption is an on-chain event on both the legacy chain and the L1, so a
prospective buyer of a legacy unit \emph{can} check whether it has been claimed --- but they must
know to look, and nothing in the token prevents the transfer. The residual risk is borne by
uninformed secondary buyers of retired-chain units during and after the window. That risk is real,
it is not mitigated by any mechanism, and its bound is the retiring exposure above.
\end{remark}

\paragraph{Reaching holders --- which now means reaching the holders who must act.} The snapshot path
removes ETH and BSC holders from this problem entirely, so the outreach target is the retiring
chains: \num{6226900.80}\flux, now measured across all eight of them
(\Cref{tab:migration-exposure}). That is a smaller and better-defined communication problem than the
one earlier drafts described, and it should be stated as smaller rather than left overstated. The
roadmap's own precedent for a parallel-asset change is announcement at least \SI{90}{\day} ahead
\srcroadmap{flux-roadmap-public.md:159}, which is \num{259200} L1 blocks \srcinferred;
$\hsnap$ would be published at announcement, i.e.\ no later than
$\hpa - \num{259200}$, together with each chain's issuer-held set $\mathcal{I}_c$ and the per-chain
dead addresses to which the Foundation transfers its own holdings, since both are needed before
$\hsnap$ \srcintent{08-answers-round-2.md, round 18}. Holders would be reached through per-asset redemption statistics
published from the snapshot, relabelled explorer token pages, exchange notices, wallet banners, and
the incumbent \Fusion endpoints answering with migration instructions for the whole window. The
single most important public act is that the Foundation immobilises its own
legacy holdings first, by transferring them to a published per-chain dead address
\srcintent{08-answers-round-2.md, round 18} --- the bulk of five EVM contracts' \num{440000000}-unit pre-mint each,
\num{2.2e9} units in total \srcinferred --- because after it, a legacy contract's remaining
\texttt{totalSupply} is public holdings only, the two disagreeing reserve-wallet lists stop
mattering, and the residual that passes to the Foundation at close is a number about third-party
holders rather than a number dominated by the Foundation's own unmoved supply.
\settlednote{Zero, gasless via \texttt{permit} (\Cref{tab:22-recommended}). Migration fee (zero recommended --- holders should not pay for a consolidation they did
not ask for), and whether \texttt{permit}/gasless migration is offered.}
\proposednote{Direct notice by zero-value calldata to holder addresses above a threshold
(approximately \num{21000} gas per holder) is this paper's proposal, with the threshold and whether
to send the notices at all unchosen; who pays for the gasless \texttt{permit} path is likewise not
recorded as decided \srcintent{08-answers-round-2.md, round 18}.}

\paragraph{Base, explicitly.} Base carries a deployed FLUX contract today
\srccode{flux-supply-tracker/\allowbreak{}index.html}{346}, added after the 2023 holder-snapshot tooling and
absent from it. The public roadmap names Base a survivor twice
\srcroadmap{flux-roadmap-public.md:97,131}; the settled decision recorded for this paper names ETH
and BSC \srcintent{03-parallel-assets.md}. Base is not BSC --- it is an Ethereum L2, and they differ
in finality, bridging architecture, audit surface and holder base. Three statements follow, and the
paper makes all three. First, under the settled decision Base is retired at $\hpa$ and its holders
redeem through the route of the preceding paragraphs. Second, \textbf{the decision is final and no later
re-addition is contemplated} \srcintent{08-answers-round-2.md, round 5} --- this closes a question
earlier drafts left open. The engineering observation that made it worth leaving open remains true
and is now only an observation: because the canonical contract is one EVM codebase, a later Base
deployment would carry near-zero marginal audit surface, so retirement is a decision about
consolidation and liquidity rather than an engineering necessity, and it should be read as such.
Third, the supersession of a published roadmap statement must be stated explicitly in the paper and
in the migration notice, because Base holders were told otherwise, and Base's measured outstanding
balance is \num{2528858.42991062}\flux
\srcmeasured{scripts/\allowbreak{}13-issuer-balances.sh}{2026-09-03} --- \SI{1.13}{\percent} of the ten-chain
total, but \SI{40.6}{\percent} of the balance actually exposed to the redemption window
(\Cref{tab:migration-exposure}) \srcinferred. Base is therefore not a rounding error in the
migration: it is most of the holder-initiated migration problem. That is a communication liability,
not an engineering one, and it is discharged by saying so.

\paragraph{Sequencing against audits and exchanges.} The order below is a hard ordering, and nothing
in it permits a public mint before the audit, before the backing is locked, or before the exchanges
hold the new address. \planned

\begin{enumerate}
\item \textbf{Audit A1} (hard gate, H7): the EVM bridge contract, the migrator
      path and \texttt{Attester\allowbreak{}Guardian} --- one codebase, published report, deployed bytecode
      matched to the audited commit.
\item \textbf{Measure} $\Lambda_c$ at $\hsnap$ and $R_0$; publish; confirm
      $R_0 \ge \Lambda$.
\item \textbf{Lock} $R_0$ into $\mathcal{S}_R$ in a public L1 transaction, with the reserve and
      redeem scripts published the same day.
\item \textbf{Deploy} ETH and BSC contracts with public mint \emph{paused}, transfers disabled
      pending $\mathsf{snapshotComplete}_c$, migrator \emph{enabled}, and the signer set published as
      four named individuals with their per-key public halves (\Cref{as:admin}), together with the
      per-chain limit $\varrho_c$ each contract enforces.
\item \textbf{Carry the snapshot} on ETH and BSC (Algorithm~\ref{alg:carry}), with exchange partners
      and ecosystem integrators coordinated in advance
      \srcintent{08-answers-round-2.md, round 7}. A custodian must know before $\hsnap$
      which of its addresses will be credited and must have decided how the credit reaches its users;
      contract-held balances are the open item of the note above. Transfers enable only when the
      carry is complete.
\item \textbf{Foundation immobilises its legacy reserves} on the retiring chains, by transfer to
      the published per-chain dead addresses \srcintent{08-answers-round-2.md, round 18}.
\item \textbf{Open} holder-initiated migration on the retiring assets and public lock-to-mint on ETH
      and BSC.
\item \textbf{Solana}: audit A2 --- the bridge program \emph{and} the SSP multisig program, the
      latter already roadmap-gated --- then steps 4--7 for Solana.
\item \textbf{Close} the window at $h^{\mathrm{mig}}_{\mathrm{close}} \approx \num{3992230}$; the
      migrator dies; the residual backing --- against the \num{6226900.80}\flux measured exposure of
      \Cref{tab:migration-exposure}, not against the ten-chain total --- passes to the Flux
      Foundation.
\end{enumerate}

\noindent Steps 3--7 land at $\hpa = \num{3466630}$, since the contract changes take effect there and
not before; steps 1 and 2 must complete earlier, and step 2 is not complete: $\Lambda_c$ is measured
for all ten legacy assets, and $R_0$ is not yet established. \textbf{Whether public mint opens under the named quorum or waits for a bonded
attester feed is settled by the same decision:} the launch configuration is 3-of-4 with autonomy to
follow, and the fallback on record is to ship 3-of-4 and take a second fork
\srcintent{08-answers-round-2.md}, so public mint opens under named signers and the attestation link
of \Cref{sec:bridging-attestation} arrives afterwards. The interval during which the pause authority
is four named individuals rather than a bonded set is therefore not an accident of sequencing; it is
the design until the second fork, and \Cref{cor:pause} is what bounds it.
\settlednote{Exchange swaps sequence \emph{before} the public window (\Cref{tab:22-recommended}). Sequencing of exchange contract swaps relative to the public window --- before, or
simultaneous. Trade-off: stranded custodial users against coordination cost with venues the project
does not control.}

\paragraph{The long tail.} Legacy contracts cannot be paused or destroyed: a plain ERC-20 has no
such function \srccode{flux-eth-bsc/\allowbreak{}flux.sol}{6--13}. They remain transferable and tradeable forever.
After the Foundation immobilises its reserves and the window closes, the remaining legacy supply is
\emph{declared} worthless but is still a token someone may sell to someone who has not read the
notice. The adversary's gain in that case is bounded by the buyer's capital and touches no bridge
state; the mitigation is the public record --- immobilised balances, the closed window, and per-chain
statistics --- which is what an explorer or an aggregator uses to mark the legacy token as migrated.

\subsection{Recommended values for the open parameters}
\label{sec:22:recommended}

The construction above leaves nineteen parameters open. Seventeen are recommended here; two are left
open because they are genuinely undetermined rather than merely unchosen, and the section says which
and why. Each was proposed by this paper and \textbf{confirmed by the author}
\srcintent{08-answers-round-2.md, round 11}, the $T_{\mathrm{react}}$ target in round 18.

\begin{table}[H]
\centering\footnotesize
\caption{Bridge parameters. Proposed by this paper and \textbf{confirmed by the author}
\srcintent{08-answers-round-2.md, round 11}, the $T_{\mathrm{react}}$ target in round 18; the two
marked \emph{open} remain so. \planned,
\srcinferred{} for the derivations except where cited.}
\label{tab:22-recommended}
\begin{tabular}{@{}llp{0.42\linewidth}@{}}
\toprule
Parameter & \textbf{Recommended} & Reason \\
\midrule
$D_{L1}$ & \textbf{\num{40} blocks} & exactly \texttt{MAX\_REORG\_LENGTH}
\srccode{fluxd/src/main.h}{74}; more adds latency and buys nothing \\
$D_{\mathrm{ETH}},D_{\mathrm{BSC}},D_{\mathrm{SOL}}$ & \textbf{each chain's finality tag} & a
confirmation \emph{count} is the BNB Bridge error class \\
$v_{\min}$ & \textbf{\num{10}\flux} & exceeds dust $+\ \varphi_{\mathrm{rel}}$ with margin; keeps a
batch entry worth its gas share \\
$\varphi_{\mathrm{mint}}$ & \textbf{0} & do not tax entry \\
$\varphi_{\mathrm{rel}}$ & \textbf{actual gas share of the batch} & cost recovery, not revenue \\
fee recipient & \textbf{Foundation operations} & not $\pool$: bridge fees must not pay node rewards \\
$T_{\mathrm{rel}}$ release cadence & \textbf{\SI{4}{\hour}} & 6 batches/day; bounded exit latency,
modest signer load \\
$\mathsf{R}_{\mathrm{hot}}$ (window) & \textbf{\num{500000}\flux} & $\approx$ 14 days at the measured
mean; covers front-loading \\
cold$\to$hot top-up & \textbf{when hot $<\ $\num{200000}\flux} & $\approx$ 6 days' cover at trigger \\
$\Delta_g$ pause guard & \textbf{\SI{48}{\hour}} & reuses the agreed upgrade timelock; one
operational constant \\
MintAuth transport & \textbf{Flux feed $+$ mandatory IPFS mirror} & IPFS is existing project
precedent \srccode{fluxd/zcutil/fetch-params.sh}{17} \\
$T_{\mathrm{att}}$ & \textbf{\num{120} blocks} (\SI{1}{\hour}) & attestation cost is per-post; hourly
is ample for a daily-scale limit \\
$T_{\mathrm{silence}}$ & \textbf{\SI{6}{\hour}} & 6 missed attestations before fail-closed \\
fail-closed on silence & \textbf{yes} & a bridge that mints while blind is the failure being designed
against \\
$T_{\mathrm{react}}$ (launch operating target) & \textbf{\SI{24}{\hour}} & the value at which
\Cref{tab:worst-case} and \Cref{prop:dominance} are evaluated; a published operational target, not a
contract constant \srcintent{08-answers-round-2.md, round 18} \\
migration fee & \textbf{0}, gasless via \texttt{permit} & holders did not ask for this consolidation \\
exchange sequencing & \textbf{before the public window} & venues need lead time; stranded custodial
users are the larger harm \\
\midrule
$(n_{\mathrm{att}}, k_{\mathrm{att}}), \beta_{\mathrm{att}}$ & \emph{open} & determined by the
slashing construction below \\
slashing construction & \emph{open} & Research-class \\
\bottomrule
\end{tabular}
\end{table}

\paragraph{Where the reserve number comes from.} $\mathsf{R}_{\mathrm{hot}}$ is the one parameter
here that can be derived rather than judged, and \Cref{tab:issuer-liability} now supplies the
operand. The measured redeeming set is \num{6226900.80}\flux{} across the eight retiring chains, and
the migration window is six months, so the \emph{mean} redemption rate over the window is
$\num{6226900.80}/180 \approx \num{34594}\flux$ per day \srcinferred. Redemption will not be uniform
--- announcement-driven flows are front-loaded, and the \SI{40.6}{\percent} of the set sitting on
Base concentrates it further --- so a hot reserve sized at the mean would be drained by the first
week. \num{500000}\flux{} is roughly fourteen days at the mean, which absorbs a first-week surge of
several times the average without a cold-storage round trip, and the \num{200000}\flux{} top-up
trigger leaves about six days of cover at the point the top-up is initiated. Both figures should be
re-derived once real redemption data exists; they are sized from the only measurement available.

\paragraph{Finality, not confirmations.} $D_{\mathrm{ETH}}$, $D_{\mathrm{BSC}}$ and
$D_{\mathrm{SOL}}$ should be expressed as each chain's own finality predicate --- Ethereum's
\texttt{finalized} block tag, BNB Chain's fast-finality tag, Solana's \texttt{finalized} commitment
--- and never as a block count. A count is a guess about reorg depth that is wrong precisely when it
matters, which is the error class the BNB Bridge incident belongs to. Using the finality gadget makes
the bridge's safety argument inherit each chain's own, which is the strongest available statement and
also the only one that stays correct when those chains change their parameters.

\paragraph{Burn-sink constructions --- superseded, and recorded for the record.} An earlier draft of
this section specified disposal constructions for the three non-EVM retiring chains: a Pact module
guard that unconditionally fails on Kadena, an unconditionally-rejecting \texttt{LogicSig} account on
Algorand, and, on Tron, the token contract's own \texttt{burn} entry point if one exists. That work is
superseded by the decision not to deploy \texttt{LegacySink} at all
(\Cref{rem:no-legacysink}): with redemption running through the incumbent for the window, no
disposal contract is deployed on any chain, EVM or otherwise, and the residual risk is carried by the
public record as that remark states. The constructions are recorded here because they remain the
right answers if a sink is ever wanted, and because the Tron case illustrates why the three chains
were not equivalent: Kadena and Algorand both admit a spending condition that is provably
unsatisfiable from the chain's own rules, whereas a Tron sink would have rested either on a
\texttt{burn} entry point whose presence this paper could not establish --- the deployed contract's
ABI is not returned by Tronscan's public contract endpoint --- or on the weaker convention that no
key exists for the zero address.

\paragraph{Custody standard for the four keys.} The recommendation is a named class rather than a
principle: keys generated and held in FIPS~140-2 Level~3 hardware or equivalent, one signer per
physical location, no signing key ever resident on a host that runs a relayer, an RPC endpoint or the
incumbent bridge service, a key-generation ceremony with published witness attestations, and a
documented rotation procedure exercised at least once before the bridge holds material value. The
last clause matters more than it looks: a rotation procedure that has never been executed is a plan,
not a control, and the first time it is attempted should not be during an incident.

\paragraph{What is left open, and why.} $(n_{\mathrm{att}}, k_{\mathrm{att}})$ and
$\beta_{\mathrm{att}}$ cannot be fixed independently of the slashing construction, because
\eqref{eq:bond-condition} is vacuous if no bond can actually be forfeited: a bond that cannot be
slashed is a deposit, and sizing it is a decision about optics rather than security. The slashing
construction is Research-class (\S\ref{sec:feasibility}), so both remain open and are not given
recommended values here. Assigning a number to $\beta_{\mathrm{att}}$ before that is settled would
be exactly the kind of plausible-looking fill this paper exists to avoid.

\subsection{What cannot be achieved}
\label{sec:bridging-open}

\begin{openproblem}[Unconditional redemption liveness against a refusing quorum]
\label{op:liveness}
Every holder can always execute \texttt{burn}: no role can prevent it, and the resulting claim is
on-chain, public and provable (H5). But the corresponding release requires $k_R$ signatures on an L1
transaction, and \Flux's inherited Bitcoin script can enforce nothing beyond ``$k_R$ of these $n_R$
keys signed'' --- it cannot enforce a rate limit, a cap, or a proof of burn. Consequently, if
$n_R - k_R + 1$ or more release signers refuse, for any reason, the holder has a provable, public,
outstanding claim and \emph{no cryptographic recourse}. At the decided $(n_R, k_R) = (4,3)$ that
number is \textbf{two}: any two of the four co-founders, by refusal, absence, incapacity or legal
order, stop every redemption to the L1 indefinitely. One refuser does not, so D1 survives against a
single uncooperative party --- but by a margin of one key, and no longer for the reason the earlier
form of this argument gave. That reason was that the Foundation holds at most $n_R - k_R$ keys and so
cannot block alone; it is not available now, because the Foundation's founders hold all four
(\Cref{as:admin}), and a correlated unavailability of two of them is a single event rather than two
independent ones. \textbf{D1 is therefore satisfied against one uncooperative party, with no margin
beyond that, and cannot be satisfied against a refusing quorum without an L1
consensus change} --- a covenant, or a bridge-vault transaction type verified by \fluxd. No such
primitive exists today. The same limitation holds for every lock-and-mint bridge whose home chain
cannot verify the foreign chain, including the incumbent, where it holds with $k = 1$.
\end{openproblem}

\noindent Two related items are open rather than solved and are recorded here so that later sections
can reference them: \Cref{op:slashing} (slashing on the L1), and the light-client construction of
\Cref{sec:bridging-alternatives}, which is closed by the header format rather than by cost. All
three are Research-class in the sense of \Cref{sec:chain-evolution}: each requires a consensus
change, not an implementation.

\subsection{Comparison with real incidents and with designs that held}
\label{sec:bridging-incidents}

Each row of \Cref{tab:bridge-incidents} is a design that the language of a bridge roadmap could be
mistaken for, together with what actually broke. The thread through the failures is not ``too
centralized''. It is \textbf{unbounded consequence}: in every failed row, once the check was defeated
--- by a bug, by a key, or by an upgrade --- the contract paid out whatever it was told to pay. The
construction's claim is correspondingly narrower than trustlessness, and defensible for that reason:
whatever is defeated, the contract issues at most $\varrho_c$ per window and at most $\bcap{c}$ ever,
visibly, and to signers whose identities are recoverable from the transaction. With the decided
parameters that ``whatever is defeated'' clause carries a number --- \num{20000000}\flux per
\SI{24}{\hour} of reaction time across the launch set of two chains (\Cref{prop:worst-case}) --- and
the third column below should be read as claiming that and nothing more. In particular, the launch
configuration does \emph{not} differ from the Ronin and Multichain rows in signer distribution; it
differs in what the contract permits after the distribution has failed. \textbf{That distinction
survives only under \Cref{rem:upgrade-voids}}: an upgradable contract without a timelock permits
whatever its next version permits, which is the Nomad row's mechanism arriving by intent rather than
by mistake.

%% longtable: nine incidents overran the fixed float's page.
{\footnotesize
\begin{longtable}{@{}p{2.0cm}p{4.6cm}p{6.4cm}@{}}
\toprule
system & what failed & how \Cref{con:bridge} differs \\
\midrule
\endfirsthead
\multicolumn{3}{@{}l}{\emph{(Table \ref{tab:bridge-incidents} continued)}}\\
\toprule
system & what failed & how \Cref{con:bridge} differs \\
\midrule
\endhead
\bottomrule
\endfoot
\textbf{Poly Network} \newline Aug 2021 &
The cross-chain manager contract could call the data contract holding the trusted keeper set, so a
crafted cross-chain message replaced the keepers with the attacker's key. \emph{The bridge's own
message path could alter its trust root.} &
No message path can touch $Q_M$, $Q_A$ or $\mathcal{G}$; only ADMIN, under $\Delta_A$, on-chain.
Mint authorisations are domain-separated by chain identifier and contract address. \\
\addlinespace
\textbf{Wormhole} \newline Feb 2022 &
Solana signature verification accepted a spoofed sysvar because a deprecated instruction-loading
function did not check the account; guardian signatures were never actually verified.
\emph{Implementation bug, not key compromise.} &
$\varrho_c$ and $\overline{M}_c$ bound what \emph{any} verification bug can issue before a pause
(\Cref{prop:bounded-mint}); the Solana Ed25519 introspection path is a named audit focus. \\
\addlinespace
\textbf{Ronin} \newline Mar 2022 &
5-of-9 validators, four run by one company plus a DAO key delegated to that same company and never
revoked; one organisation was socially engineered; undetected for six days. \emph{Threshold nominal,
concentration real, no monitoring.} &
\textbf{Not by distribution.} At launch all four keys sit with one organisation's founders
(\Cref{as:admin}), which is the Ronin concentration and this paper says so. This row is also the
answer to why \Cref{prop:worst-case} models a 3-of-4 compromise at all: a 5-of-9 quorum was taken
whole. What differs: every signature is attributable to a named key; the pause is 1-of-4, so a
3-key compromise does not control it (\Cref{cor:pause}); and issuance after a full compromise is
bounded to \num{20000000}\flux per \SI{24}{\hour} of detection latency across two chains
(\Cref{prop:worst-case}) rather than to the whole reserve. Ronin's six undetected days exceed
$\Delta_A$: the rate limit would have held the loss to $2 \times \num{20000000}\flux$ for the
\qty{48}{\hour} until a queued upgrade landed, and the contract-enforced ceiling thereafter is
$\sum_c \bcap{c}$ (\Cref{rem:global-cap}), not the rate limit's number. \\
\addlinespace
\textbf{Harmony Horizon} \newline Jun 2022 &
2-of-5 multisig; two keys compromised. \emph{Threshold too low for the value at risk.} &
$(n,k) = (4,3)$ decided (\Cref{tab:bridge-decided}): a higher threshold ratio than Harmony's on a
smaller set, so two independent key compromises are survivable and three are not. The loss is
bounded by \Cref{prop:worst-case} or $\mathsf{R}_{\mathrm{hot}}$ regardless. \\
\addlinespace
\textbf{Nomad} \newline Aug 2022 &
An upgrade initialised the trusted root to zero, so every unproven message verified and anyone could
copy the exploit. \emph{An in-place upgrade regressed a security invariant.} &
\textbf{This row is answered only by the timelock.} The construction specifies no upgradeability (D3), with
replacement only by holder-initiated migration to a separately audited contract; the decided
configuration keeps an upgrade path under the same 3-of-4
\srcintent{08-answers-round-2.md, round 7}. Under D3 the Nomad mechanism is unavailable; under the
decision it is available to three keys, and the \qty{48}{\hour} timelock of \Cref{rem:timelock-decided} is what
makes it slow and public instead. \\
\addlinespace
\textbf{BNB Bridge} \newline Oct 2022 &
A flaw in the IAVL Merkle-proof verification used by the light client let a forged proof pass;
validators halted the chain to stop it. \emph{Even a light client is only as good as its proof
verifier, and there was no rate limit.} &
$\varrho_c$ per window makes an equivalent forgery worth $\varrho_c$ rather than the cap, and the
pause is a contract action rather than a chain halt. \\
\addlinespace
\textbf{Multichain} \newline Jul 2023 &
MPC ``nodes'' whose keys and servers were under one person's sole control; his arrest left the
bridge inoperable, then drained. \emph{A threshold with one holder is one key.} &
Per-signer keys held by four separately named individuals, not shares of one aggregate key, so the
contract sees which three signed and a single holder cannot produce a valid authorisation. The
concentration risk is not removed --- see the Ronin row --- and the incumbent's 32-key model is the
same failure class, which this paper says rather than implies. \\
\midrule
\textbf{Circle CCTP} \citep{cctp} \newline held &
Issuer-attested burn-and-mint; no liquidity pool to steal; nonce-bound messages. Trust is the
issuer, who is already trusted for the whole supply. &
Adopted as the model. Adds explicit $k$-of-$n$ attribution and contract-level bounds, because
\Flux's attesters are \emph{not} already trusted for the whole supply. \\
\addlinespace
\textbf{Optimism / Arbitrum} \newline held &
L1 verifies L2; challenge window; guardian pause; security-council timelock. &
Cross-verification is unavailable (\Cref{sec:bridging-alternatives}); the guardian, timelock and
no-silent-upgrade pattern is adopted. \\
\addlinespace
\textbf{IBC} \citep{cosmosibc} \newline held &
Light clients in both directions on fast-finality chains. &
Not applicable without a node-set commitment in the \PoN header --- the constraint of
\Cref{sec:bridging-alternatives}, not a preference. \\
\caption{Bridge incidents and designs that held, against \Cref{con:bridge}. Loss figures are omitted
deliberately: they are widely reported, they vary by source and by the price used, and the argument
does not depend on them. What the argument depends on is the failure mechanism in the middle column.
\srcinferred, from published incident post-mortems; see \Cref{sec:bridging-alternatives} for the
design references.}
\label{tab:bridge-incidents}\\
\end{longtable}
}

\paragraph{Feasibility, in one paragraph.} The canonical EVM contract, the migrator, and the
off-chain attestation feed with its on-chain guardian are \emph{engineering}: standard Solidity and
standard verification, where the novelty is discipline rather than cryptography. The L1 reserve as a
published P2SH set, the memo-carried ledgers, and the anyone-verifiable $\varsigma$ of
Algorithm~\ref{alg:verify} are \emph{extensions} of primitives the chain already has and the project
already uses. Per-node, L1-registered attester keys are engineering, but they replace a fleet-shared
derivation and their registration transaction is new design. Two items are Research:
\Cref{op:slashing}, and a \Flux light client on an EVM chain, which is gated by the header format
rather than by cost. One item --- \Cref{op:liveness} --- is not achievable at all without a consensus
change. A third is not classified here at all: the autonomous replacement for the 3-of-4 quorum has
no construction in the corpus, so it is neither engineering nor Research until R1--R5 of
\Cref{sec:bridging-assumption} have a candidate answer. \Cref{sec:chain-evolution} carries the header
question and \Cref{sec:attested-execution} carries the attester-key question; this section carries
the rest.

%% Drain this section's float queue before the next begins.
\FloatBarrier
